\chapter{Całka Kurzweila--Stieltjesa}
\providecommand{\atlasref}[2]{#2}
\label{ma1:ch:calka-kurzweila-stieltjesa}
\begingroup
\providecommand{\osc}{\operatorname{osc}}
\providecommand{\Var}{\operatorname{Var}}
\providecommand{\mesh}{\operatorname{mesh}}
\providecommand{\supp}{\operatorname{supp}}
\providecommand{\AC}{\mathrm{AC}}
\providecommand{\BV}{\mathrm{BV}}
\providecommand{\HK}{\mathrm{KH}}
\providecommand{\ind}[1]{\chi_{#1}}
\providecommand{\Dn}{\underline D}
\providecommand{\Du}{\overline D}
\providecommand{\KS}{\mathrm{KS}}
Przyjmujemy oznaczenia i konwencje z rozdziału \ref{ma1:ch:kurzweil-henstock}.
W szczególności $V_u^v(g)$ oznacza wahanie całkowite
$g$ na $[u,v]$, a $\BV[a,b]$ klasę funkcji o ograniczonym
wahaniu. Funkcję $g:[a,b]\to\R$ nazywamy regulowaną,
jeżeli istnieją skończone granice $g(t-)$ dla $t\in(a,b]$
oraz $g(t+)$ dla $t\in[a,b)$.
Funkcja schodkowa jest stała na otwartych komórkach
pewnego skończonego podziału, przy czym jej wartości
w węzłach mogą być dowolne. Równoważność regulowania
i jednostajnej aproksymowalności takimi funkcjami
wykażemy poniżej.
\section{Całka Kurzweila--Stieltjesa}
\label{sec:kurzweil-stieltjes}
Całka Kurzweila--Stieltjesa (KS) łączy przyrosty funkcji całkującej z cechowaniami znanymi z całki Kurzweila--Henstocka. Jej przewaga nad całką Riemanna--Stieltjesa ujawnia się przede wszystkim przy wspólnych nieciągłościach funkcji podcałkowej i całkującej. W tej sekcji, oprócz kryteriów istnienia, wyprowadzimy wzory obliczeniowe i oszacowania błędu przybliżeń.

\subsection*{Definicja, jednoznaczność i podstawowe oszacowania}
\begin{definicja}[Całka Kurzweila--Stieltjesa]
Niech $a<b$ i niech $f,\alpha:[a,b]\to\R$ będą funkcjami o skończonych wartościach. \emph{Cechowaniem} nazywamy funkcję $\delta:[a,b]\to(0,\infty)$. Podział znaczony $\mathcal P=\{(\xi_i,[x_{i-1},x_i])\}_{i=1}^n$, gdzie $a=x_0<\cdots<x_n=b$ i $\xi_i\in[x_{i-1},x_i]$, jest \emph{$\delta$-drobny}, jeżeli $[x_{i-1},x_i]\subset(\xi_i-\delta(\xi_i),\xi_i+\delta(\xi_i))$ dla każdego $i$. Definiujemy
\[
 S_\alpha(f,\mathcal P):=\sum_{i=1}^n f(\xi_i)\bigl(\alpha(x_i)-\alpha(x_{i-1})\bigr).
\]
Liczba $I\in\R$ jest całką KS z $f$ względem $\alpha$, jeżeli dla każdego $\varepsilon>0$ istnieje cechowanie $\delta$ takie, że $|S_\alpha(f,\mathcal P)-I|<\varepsilon$ dla każdego $\delta$-drobnego podziału znaczonego $\mathcal P$. Piszemy wówczas $I=(\KS)\int_a^b f\,\mathrm d\alpha$, a klasę takich funkcji $f$ oznaczamy przez $\mathrm{KS}(\alpha)$. Na przedziale jednopunktowym przyjmujemy całkę zero.
\end{definicja}
Istnienie podziału $\delta$-drobnego dla każdego cechowania wynika z lematu Cousina (lemat \ref{ma1:kh:lem:kh-cousin}). Gdy $\alpha(x)=x$, definicja jest dokładnie definicją całki Kurzweila--Henstocka. W szczególności nie wolno utożsamiać dowolnej całki KS z całką względem miary Stieltjesa: potrzebne są dodatkowe założenia.

\begin{lemat}[Wahanie a regulowanie]\label{it:lem:ks-bv-regulated}
Każda funkcja $g\in\BV[a,b]$ jest regulowana.
\end{lemat}
\begin{proof}
Niech $V(x)=V_a^x(g)$, przy czym $V(a)=0$. Addytywność wahania (lemat \ref{it:lem:rs-variation}) daje dla $u<v$ nierówność $V(v)-V(u)\ge|g(v)-g(u)|$. Funkcje $(V+g-g(a))/2$ i $(V-g+g(a))/2$ są zatem niemalejące i ograniczone, gdyż $0\le V\le V(b)$ oraz $|g-g(a)|\le V(b)$. Każda ograniczona funkcja niemalejąca ma skończone granice jednostronne, równe odpowiednim kresom jej wartości po danej stronie punktu. Różnica tych dwóch funkcji plus stała $g(a)$ jest równa $g$, a więc także ma skończone granice jednostronne.
\end{proof}

\begin{twierdzenie}[Kryterium Cauchy'ego, liniowość i zgodność]
\label{tw:ks-cauchy}
Dla funkcji $f,\alpha:[a,b]\to\R$ o skończonych wartościach całka $(\KS)\int_a^b f\,\mathrm d\alpha$ istnieje wtedy i tylko wtedy, gdy dla każdego $\varepsilon>0$ istnieje cechowanie $\delta$ takie, że
$|S_\alpha(f,\mathcal P)-S_\alpha(f,\mathcal Q)|<\varepsilon$ dla dowolnych dwóch $\delta$-drobnych podziałów znaczonych $\mathcal P,\mathcal Q$. Całka jest jednoznaczna i liniowa osobno w $f$ i $\alpha$, o ile rozważane całki istnieją. Jeżeli istnieje właściwa całka Riemanna--Stieltjesa $\int_a^b f\,\mathrm d\alpha$, to istnieje również całka KS i ma tę samą wartość.
\end{twierdzenie}
\begin{proof}
Dla $a=b$ pusty podział ma sumę zero, więc tezy są bezpośrednie. Dalej zakładamy $a<b$. Jeżeli całka istnieje i $\varepsilon>0$, wybieramy cechowanie z definicji całki dla błędu $\varepsilon/2$; różnica dwóch drobnych sum jest wtedy mniejsza od $\varepsilon$. Odwrotnie, dla każdego $m\ge1$ wybieramy cechowanie $\delta_m$ zapewniające różnice sum mniejsze niż $2^{-m}$ i zastępujemy je przez $\min(\delta_1,\ldots,\delta_m)$. Z lematu Cousina wybieramy $\delta_m$-drobne podziały $\mathcal P_m$. Gdy $k\ge m$, obie sumy $S_\alpha(f,\mathcal P_k)$ i $S_\alpha(f,\mathcal P_m)$ są porównywalne przez cechowanie $\delta_m$, toteż tworzą ciąg Cauchy'ego. Niech $I$ będzie jego granicą. Dla dowolnego $\delta_m$-drobnego podziału $\mathcal P$ porównujemy jego sumę z $\mathcal P_k$ i przechodzimy z $k$ do nieskończoności; otrzymujemy $|S_\alpha(f,\mathcal P)-I|\le2^{-m}$. Dla dowolnego $\varepsilon>0$ wybieramy $m$ z $2^{-m}<\varepsilon$, więc $\delta_m$ spełnia definicję całki. Jeżeli $I,J$ spełniają definicję całki, dla
$\varepsilon>0$ bierzemy minimum cechowań dla obu wartości
z błędem $\varepsilon/2$. Lemat Cousina daje podział drobny
dla tego minimum, a nierówność trójkąta daje $|I-J|<\varepsilon$.
Dowolność $\varepsilon$ daje $I=J$.
Dla kombinacji liniowej $s f_1+t f_2$ przy ustalonym
integratorze wybieramy cechowania dla obu całek z błędem
$\varepsilon/(1+|s|+|t|)$ i bierzemy ich minimum.
Liniowość sum daje błąd mniejszy od $\varepsilon$ względem
odpowiedniej kombinacji wartości całek. Ten sam argument
przy ustalonej funkcji podcałkowej i kombinacji integratorów
dowodzi liniowości względem integratora. Dla całki Riemanna--Stieltjesa wybieramy cechowanie stałe równe połowie dopuszczalnej średnicy podziału, gdyż komórka drobna ma długość mniejszą od dwukrotności promienia cechowania.
\end{proof}

\begin{twierdzenie}[Lokalność i addytywność]\label{it:thm:ks-locality}
Niech $f,g:[a,b]\to\R$ mają skończone wartości. Jeśli istnieje
całka KS na $[a,b]$, istnieje na każdym podprzedziale.
Dla $a<c<b$ całkowalność na $[a,b]$ jest równoważna
całkowalności na obu $[a,c]$ i $[c,b]$; wtedy
\[
 (\KS)\int_a^b f\,\mathrm dg=
 (\KS)\int_a^c f\,\mathrm dg+(\KS)\int_c^b f\,\mathrm dg.
\]
Na przedziale jednopunktowym przyjmujemy całkę zero.
\end{twierdzenie}
\begin{proof}
Jeśli istnieje całka na całości, wybieramy cechowanie z kryterium
Cauchy'ego (twierdzenie \ref{tw:ks-cauchy}). Dwa drobne podziały
$[a,c]$ uzupełniamy tym samym drobnym podziałem $[c,b]$,
istniejącym dzięki lematowi Cousina \ref{ma1:kh:lem:kh-cousin}.
Różnica pełnych sum jest równa różnicy sum na $[a,c]$.
Kryterium Cauchy'ego daje całkowalność na $[a,c]$; zamiana ról
przedziałów daje ją na $[c,b]$, a powtórzenie argumentu na
podprzedziałach daje lokalność.

Odwrotnie, wybieramy cechowania na obu częściach dla błędu
$\varepsilon/2$. Poza $c$ ograniczamy je dodatkowo przez
$|t-c|/2$, a w $c$ bierzemy minimum obu wartości.
Każda komórka zawierająca $c$ ma znacznik $c$,
ponieważ dla $t\ne c$ promień jest mniejszy niż $|t-c|$.
Komórkę zawierającą $c$ we wnętrzu rozcinamy w $c$,
zachowując znacznik w obu częściach. Jeżeli $c$ jest
już węzłem, rozcięcie nie jest potrzebne.
Wkład rozcinanej komórki nie zmienia się: jej wkład nie zmienia
się, ponieważ $g(v)-g(u)=g(c)-g(u)+g(v)-g(c)$.
Otrzymane podziały obu części są drobne dla odpowiednich
cechowań, zatem pełna suma różni się od sumy całek lokalnych
o mniej niż $\varepsilon$. Jednoznaczność kończy dowód.
\end{proof}

\begin{lemat}[Saksa--Henstocka dla całki KS]\label{it:lem:ks-saks-henstock}
Niech $f$ będzie KS-całkowalna względem $g$,
$F(x)=(\KS)\int_a^x f\,\mathrm dg$, a cechowanie $\delta$
zapewnia błąd mniejszy od $\varepsilon$ na $[a,b]$.
Dla każdego skończonego $\delta$-drobnego podziału
częściowego $\{(t_j,[u_j,v_j])\}_{j=1}^m$ o rozłącznych
wnętrzach zachodzi
\[
 \sum_j|f(t_j)(g(v_j)-g(u_j))-F(v_j)+F(u_j)|
 \le 2\varepsilon.
\]
\end{lemat}
\begin{proof}
Niech $e_j=f(t_j)(g(v_j)-g(u_j))-F(v_j)+F(u_j)$.
Domknięcia luk o dodatniej długości między komórkami
oraz przy końcach przedziału tworzą co najwyżej $m+1$
przedziałów. Dla dowolnego $\eta>0$ na każdej luce
wybieramy podział drobny względem minimum $\delta$ i
cechowania dla całki lokalnej z błędem $\eta/(m+1)$.
Lokalność z twierdzenia \ref{it:thm:ks-locality} oraz
lemat Cousina zapewniają możliwość takiego wyboru.
Połączenie podziałów luk z podziałem częściowym jest
pełnym podziałem $\delta$-drobnym; przyrosty $F$ teleskopują,
a błędy w lukach mają łącznie moduł mniejszy od $\eta$.
Otrzymujemy $|\sum_j e_j|\le\varepsilon+\eta$,
a po przejściu z $\eta\downarrow0$ mamy
$|\sum_j e_j|\le\varepsilon$.
Ten sam argument stosujemy do dowolnego podzbioru komórek,
w szczególności osobno do tych z $e_j\ge0$ i $e_j<0$.
Dodanie obu oszacowań daje sumę modułów nie większą niż
$2\varepsilon$.
\end{proof}

\begin{twierdzenie}[Charakteryzacja funkcji regulowanych]
Funkcja $f:[a,b]\to\R$ jest regulowana wtedy i tylko wtedy,
gdy jest jednostajną granicą funkcji schodkowych.
W szczególności jest ograniczona, borelowska i całkowalna
w sensie Riemanna.
\end{twierdzenie}
\begin{proof}
Dla $a=b$ każda funkcja jest schodkowa i wszystkie tezy
są bezpośrednie. Zakładamy dalej $a<b$.
Każda funkcja schodkowa ma granice jednostronne. Jeśli $f_n$ mają granice jednostronne i $f_n\to f$ jednostajnie, to dla $\varepsilon>0$ wybieramy $n$ z $\|f-f_n\|_\infty<\varepsilon/3$. Przy ustalonej dopuszczalnej stronie punktu $t$ wybieramy $\rho>0$ tak, aby $|f_n(u)-f_n(v)|<\varepsilon/3$ dla $u,v$ leżących po tej samej stronie $t$ w odległości mniejszej niż $\rho$. Wtedy
\[
 |f(u)-f(v)|\le 2\|f-f_n\|_\infty+|f_n(u)-f_n(v)|<\varepsilon.
\]
Jednostronny warunek Cauchy'ego i zupełność $\R$ dają istnienie skończonej granicy jednostronnej $f$. W przeciwnym kierunku ustalmy $\varepsilon>0$. W każdym punkcie $t$ mamy jednostronne otoczenia, na których oscylacja $f$ jest mniejsza niż $\varepsilon$. Wybieramy cechowanie $\delta(t)>0$, którego lewe i prawe półotoczenia bez punktu $t$ mają oscylację mniejszą niż $\varepsilon$ (po przecięciu z dziedziną). Lemat \ref{ma1:kh:lem:kh-cousin} dostarcza skończony $\delta$-drobny podział. Wstawiamy wszystkie jego znaczniki do zbioru węzłów i usuwamy powtórzenia. Każdy nowy otwarty przedział jest zawarty w jednym z wymienionych półotoczeń znacznika pierwotnej komórki, więc jego oscylacja jest mniejsza niż $\varepsilon$. Na tych przedziałach podstawiamy stałe, a w węzłach zachowujemy dokładne wartości $f$. Powstała funkcja schodkowa przybliża $f$ jednostajnie z błędem co najwyżej $\varepsilon$.
W szczególności funkcja regulowana jest ograniczona i borelowska jako jednostajna granica funkcji schodkowych. Każda funkcja schodkowa jest całkowalna w sensie Riemanna, więc twierdzenie \ref{ma1:thm:riemann-uniform-limit} daje całkowalność Riemanna każdej funkcji regulowanej.
\end{proof}

\begin{lemat}[Oszacowanie przez wahanie]
\label{lem:ks-var}
Jeżeli $\alpha\in\BV[a,b]$ i $f:[a,b]\to\R$ jest ograniczona, to dla każdego podziału znaczonego $\mathcal P$ zachodzi $|S_\alpha(f,\mathcal P)|\le\|f\|_\infty V_a^b(\alpha)$. Jeżeli całka KS istnieje, to
\[
 \left|(\KS)\int_a^b f\,\mathrm d\alpha\right|\le\|f\|_\infty V_a^b(\alpha).
\]
\end{lemat}
\begin{proof}
Dla dowolnego podziału znaczonego
\[
 |S_\alpha(f,\mathcal P)|
 \le\sum_i|f(\xi_i)|\,|\alpha(x_i)-\alpha(x_{i-1})|
 \le\|f\|_\infty V_a^b(\alpha).
\]
Jeżeli istnieje całka $I$, dla każdego $\varepsilon>0$
wybieramy cechowanie z jej definicji, a z lematu Cousina
podział drobny dla tego cechowania (dla $a=b$ używamy
pustego podziału). Wówczas
$|I|\le |I-S_\alpha(f,\mathcal P)|+|S_\alpha(f,\mathcal P)|
<\varepsilon+\|f\|_\infty V_a^b(\alpha)$.
Dowolność $\varepsilon$ daje tezę.
\end{proof}

\subsection*{Skoki, wartości w punktach i funkcje schodkowe}
Dla regulowanej $g:[a,b]\to\R$ stosujemy oznaczenia $\Delta^-g(t)=g(t)-g(t-)$ dla $t\in(a,b]$, $\Delta^+g(t)=g(t+)-g(t)$ dla $t\in[a,b)$ i $\Delta g(t)=g(t+)-g(t-)$ dla $t\in(a,b)$. Dla $a<b$ na końcach przyjmujemy $\Delta^-g(a)=\Delta^+g(b)=0$ oraz $\Delta g(a)=\Delta^+g(a)$, $\Delta g(b)=\Delta^-g(b)$. Jeżeli $a=b$, wszystkie trzy skoki przyjmujemy równe zeru; w tym przypadku nie używamy granic jednostronnych. Rozróżnienie skoku lewostronnego i prawostronnego jest istotne: $g(t)$ nie musi równać się żadnej granicy jednostronnej.

\begin{lemat}[Punktowy wzór wykonawczy]
\label{lem:ks-point}
Niech $g:[a,b]\to\R$ będzie regulowana i $c\in[a,b]$. Dla dowolnej liczby $q\in\R$ zachodzi
\[
 (\KS)\int_a^b q\ind{\{c\}}(t)\,\mathrm dg(t)=q\Delta g(c).
\]
W szczególności wartość funkcji podcałkowej w punkcie o $\Delta g(c)\ne0$ wpływa na całkę; dwa przeciwne skoki jednostronne mogą się natomiast znosić.
\end{lemat}
\begin{proof}
Dla $a=b$ teza wynika z konwencji zerowej całki i zerowych skoków. Ustalmy dokładność $\eta>0$. Jeżeli $q=0$, teza wynika natychmiast z zerowych sum. W przeciwnym razie kładziemy $\varepsilon=\eta/(3|q|)$. Rozpatrzmy $a<c<b$. Dla $t\ne c$ wybieramy $\delta(t)<|t-c|/2$; wartość $\delta(c)>0$ dobieramy tak małą, aby $|g(u)-g(c-)|<\varepsilon$ dla $c-\delta(c)<u<c$ i $|g(v)-g(c+)|<\varepsilon$ dla $c<v<c+\delta(c)$. W podziale drobnym każda komórka, której wnętrze zawiera $c$, musi mieć znacznik $c$. Jeżeli $c$ jest węzłem podziału, jedynie dwie przyległe do niego komórki mogą mieć znacznik $c$; obie rzeczywiście muszą go mieć przy dodatniej długości. Niezerowe składniki sumy teleskopują więc do $q(g(v)-g(u))$, gdzie $u<c<v$: jeśli $c$ nie jest węzłem, są to końce jednej komórki, a w przeciwnym razie zewnętrzne końce dwóch sąsiednich komórek. Ich odległości od $c$ są mniejsze niż $\delta(c)$. Wartość sumy różni się od $q(g(c+)-g(c-))$ o mniej niż $2|q|\varepsilon<\eta$. Dla $c=a$ wybieramy $\delta(a)$ z $|g(v)-g(a+)|<\varepsilon$ dla $a<v<a+\delta(a)$; pierwsza komórka musi mieć znacznik $a$ i suma jest równa $q(g(v)-g(a))$. Błąd względem $q\Delta^+g(a)$ jest mniejszy niż $|q|\varepsilon<\eta$. Dla $c=b$ wybieramy analogicznie $\delta(b)$ z $|g(u)-g(b-)|<\varepsilon$ dla $b-\delta(b)<u<b$; ostatnia komórka ma znacznik $b$, a suma $q(g(b)-g(u))$ różni się od $q\Delta^-g(b)$ o mniej niż $\eta$. To dowodzi istnienia całki we wszystkich przypadkach.
\end{proof}

\begin{lemat}[Funkcje charakterystyczne przedziałów]
\label{lem:ks-indicators}
Jeśli $g:[a,b]\to\R$ jest regulowana, to dla $a\le u<v\le b$ istnieją całki KS z funkcji charakterystycznych $[u,v]$, $(u,v)$, $[u,v)$ i $(u,v]$. Przykładowo dla $a<u<v<b$ otrzymujemy
\[
 \begin{aligned}
 (\KS)\int_a^b\ind{(u,v)}\,\mathrm dg&=g(v-)-g(u+),\\
 (\KS)\int_a^b\ind{[u,v]}\,\mathrm dg&=g(v+)-g(u-),\\
 (\KS)\int_a^b\ind{[u,v)}\,\mathrm dg&=g(v-)-g(u-),\\
 (\KS)\int_a^b\ind{(u,v]}\,\mathrm dg&=g(v+)-g(u+).
 \end{aligned}
\]
Gdy $u=a$ lub $v=b$, symbole $g(a-)$ i $g(b+)$ we wzorach zastępujemy odpowiednio $g(a)$ i $g(b)$.
\end{lemat}
\begin{proof}
Dla funkcji $\ind{(u,v)}$ wybieramy cechowanie tak, aby komórki o znacznikach różnych od $u,v$ nie zawierały tych punktów: dla takiego znacznika $x$ wystarczy $\delta(x)<\min(|x-u|,|x-v|)$. Dla znaczników $u,v$ wybieramy $\delta(u),\delta(v)<(v-u)/3$. Komórki zawierające $u$ lub $v$ muszą być oznaczone odpowiednio tymi punktami, więc wnoszą zero. Komórki ze znacznikami wewnątrz $(u,v)$ tworzą niepusty blok pokrywający $[s,r]$ z $u<s<r<v$, $s-u<\delta(u)$ i $v-r<\delta(v)$. Rzeczywiście, skrajne komórki zawierające $u$ i $v$ sięgają do wnętrza $(u,v)$, lecz ich zasięgi nie mogą się zetknąć z powodu ograniczenia promieni. Wszystkie komórki między nimi mają znaczniki wewnętrzne. Ich przyrosty teleskopują do $g(r)-g(s)$. Dla zadanej dokładności zmniejszamy $\delta(u),\delta(v)$ tak, aby z istnienia granic jednostronnych $g$ wynikało $|g(s)-g(u+)|<\varepsilon/2$ oraz $|g(r)-g(v-)|<\varepsilon/2$. Daje to pierwszy wzór także przy $u=a$ lub $v=b$. Następnie korzystamy z tożsamości funkcji charakterystycznych, np. $\ind{[u,v]}=\ind{(u,v)}+\ind{\{u\}}+\ind{\{v\}}$, oraz z lematu~\ref{lem:ks-point}. Pozostałe wzory uzyskujemy analogicznie.
\end{proof}

\begin{twierdzenie}[Całkowanie względem funkcji schodkowej]
\label{tw:ks-step}
Niech $g:[a,b]\to\R$ będzie funkcją schodkową, mającą skończony zbiór punktów nieciągłości $D$, i niech $f:[a,b]\to\R$ będzie dowolną funkcją o skończonych wartościach. Wtedy całka KS istnieje i
\[
 (\KS)\int_a^b f\,\mathrm dg
 =f(a)\Delta^+g(a)+\sum_{c\in D\cap(a,b)}f(c)\Delta g(c)+f(b)\Delta^-g(b).
\]
\end{twierdzenie}
\begin{proof}
Dla $a=b$ teza wynika z konwencji całki zerowej. Niech $a=c_0<c_1<\cdots<c_m=b$ będzie podziałem, na którego otwartych komórkach $g$ jest stała. Wówczas
\[
 g(t)=g(a)+\sum_{j=1}^m\Delta^-g(c_j)\chi_{[c_j,b]}(t)
 +\sum_{j=0}^{m-1}\Delta^+g(c_j)\chi_{(c_j,b]}(t).
\]
Aby to sprawdzić, zaczynamy w $a$, gdzie obie sumy są zerowe. Przejście z węzła $c_j$ do następnej otwartej komórki dodaje dokładnie $\Delta^+g(c_j)$, a przejście z tej komórki do $c_{j+1}$ dodaje dokładnie $\Delta^-g(c_{j+1})$. Indukcyjnie prawa strona odtwarza wartości $g$ zarówno we wszystkich węzłach, jak i na wszystkich otwartych komórkach. Węzły, w których $g$ jest ciągła, mają oba skoki zerowe. Dla $c\in(a,b]$ i integratora $\ind{[c,b]}$ jedyna komórka z niezerowym przyrostem ma końce $u<c\le v$. Wybierając dla $t\ne c$ cechowanie $\delta(t)<|t-c|/2$, wymuszamy znacznik $c$ na tej komórce, a suma wynosi dokładnie $f(c)$. Dla $c\in[a,b)$ i integratora $\ind{(c,b]}$ jedyna taka komórka ma końce $u\le c<v$; identyczny wybór cechowania daje $f(c)$. Przy $c=a$ pierwszy integrator jest stały, a przy $c=b$ drugi jest zerowy, więc nie wnoszą wkładu. Liniowość względem integratora daje wskazany wzór. Nie potrzebujemy ograniczoności ani regularności $f$, ponieważ występuje skończenie wiele wartości $f(c)$.
\end{proof}

\begin{przyklad}[Wspólny skok i nieistnienie zwykłej całki Stieltjesa]
Niech $a<c<b$ i $H=\ind{[c,b]}$. Wtedy dla każdej funkcji $f$ o skończonych wartościach zachodzi $(\KS)\int_a^b f\,\mathrm dH=f(c)$. Stąd $(\KS)\int_a^b H\,\mathrm dH=1$, a $(\KS)\int_a^b\ind{(c,b]}\,\mathrm dH=0$. Pierwsza z tych całek nie istnieje jako zwykła całka Riemanna--Stieltjesa, ponieważ swobodny wybór znacznika na komórce przecinającej $c$ daje różne sumy. Całka KS rozstrzyga tę niejednoznaczność za pomocą cechowania.
\end{przyklad}

\begin{przyklad}[Punktowy impuls integratora]
Niech $g=\ind{\{c\}}$ dla $a<c<b$. Wtedy $\Delta^-g(c)=1$ i $\Delta^+g(c)=-1$, zatem $\Delta g(c)=0$. Z twierdzenia \ref{tw:ks-step} dla każdej funkcji
$f:[a,b]\to\R$ o skończonych wartościach otrzymujemy
$(\KS)\int_a^b f\,\mathrm dg=0$; nie jest tu potrzebna
regulowalność ani ograniczoność $f$. Oznacza to, że dwie przeciwne zmiany funkcji całkującej, następujące po obu stronach tego samego punktu, znoszą się. Wartość $g(c)$ może więc nie wpływać na całkę, mimo że oddzielne skoki jednostronne są niezerowe.
\end{przyklad}

\subsection*{Dwa podstawowe twierdzenia o istnieniu}
\begin{twierdzenie}[Funkcja regulowana i całkująca o wahaniu ograniczonym]
\label{tw:ks-reg-bv}
Jeśli $f:[a,b]\to\R$ jest regulowana, a $g\in\BV[a,b]$, to całka $(\KS)\int_a^b f\,\mathrm dg$ istnieje. Ponadto
$|(\KS)\int_a^b f\,\mathrm dg|\le\|f\|_\infty V_a^b(g)$.
\end{twierdzenie}
\begin{proof}
Z charakteryzacji funkcji regulowanych wybieramy funkcje schodkowe $f_n$ z $\|f_n-f\|_\infty\to0$. Funkcja $f$ jest ograniczona: dla dowolnego $n$ z $\|f-f_n\|_\infty<1$ mamy $\|f\|_\infty\le1+\|f_n\|_\infty<\infty$. Każda całka $I_n=(\KS)\int_a^b f_n\,\mathrm dg$ istnieje na mocy lematów~\ref{lem:ks-indicators} i~\ref{lem:ks-point}, ponieważ $f_n$ jest skończoną kombinacją funkcji charakterystycznych przedziałów i punktów, a $g$ jest regulowana na mocy rozkładu Jordana z lematu \ref{ma1:kh:lem:variation-additivity}. Z lematu~\ref{lem:ks-var} i liniowości wynika $|I_n-I_m|\le\|f_n-f_m\|_\infty V_a^b(g)$. Niech $I=\lim_n I_n$. Dla dowolnego podziału znaczonego $\mathcal P$ mamy
\[
 |S_g(f,\mathcal P)-I|\le
 \|f-f_n\|_\infty V_a^b(g)+|S_g(f_n,\mathcal P)-I_n|+|I_n-I|.
\]
Dla ustalonego $\varepsilon>0$ wybieramy $n$ tak duże, aby $\|f-f_n\|_\infty V_a^b(g)<\varepsilon/3$ i $|I_n-I|<\varepsilon/3$. Z definicji całki $I_n$ wybieramy cechowanie zapewniające $|S_g(f_n,\mathcal P)-I_n|<\varepsilon/3$ dla każdego drobnego podziału $\mathcal P$. Powyższa nierówność daje wtedy $|S_g(f,\mathcal P)-I|<\varepsilon$, co dowodzi istnienia całki. Ostatnia nierówność wynika z lematu~\ref{lem:ks-var} po przejściu do granicy.
\end{proof}

\begin{lemat}[Nierówność sumowania przez części]
\label{lem:ks-summation}
Niech $f\in\BV[a,b]$ i niech $g:[a,b]\to\R$ będzie ograniczona. Dla każdego podziału znaczonego $\mathcal P$ zachodzi
\[
 |S_g(f,\mathcal P)|\le
 \bigl(2\|f\|_\infty+V_a^b(f)\bigr)\|g\|_\infty.
\]
\end{lemat}
\begin{proof}
Oznaczmy $f_i=f(\xi_i)$, $g_i=g(x_i)$. Przekształcenie skończonej sumy daje
$S_g(f,\mathcal P)=f_ng_n-f_1g_0+\sum_{i=1}^{n-1}(f_i-f_{i+1})g_i$. Ponieważ $\xi_i\in[x_{i-1},x_i]$ oraz
$\xi_{i+1}\in[x_i,x_{i+1}]$, mamy $\xi_i\le\xi_{i+1}$.
Usuwamy powtarzające się znaczniki (odpowiadające im
przyrosty są zerowe) i dołączamy końce $a,b$, jeżeli
jeszcze nie występują. Otrzymujemy podział $[a,b]$,
którego suma modułów przyrostów $f$ jest nie mniejsza
niż $\sum_{i=1}^{n-1}|f_i-f_{i+1}|$. Z definicji wahania
ta ostatnia suma nie przekracza $V_a^b(f)$. Oszacowanie wynika z nierówności trójkąta.
\end{proof}

\begin{twierdzenie}[Całkująca regulowana i funkcja podcałkowa o wahaniu ograniczonym]
\label{tw:ks-bv-reg}
Jeśli $f\in\BV[a,b]$ i $g:[a,b]\to\R$ jest regulowana, to całka $(\KS)\int_a^b f\,\mathrm dg$ istnieje. Jeżeli $g,h$ są regulowane, to
\[
 \left|(\KS)\int_a^b f\,\mathrm dg-(\KS)\int_a^b f\,\mathrm dh\right|
 \le\bigl(2\|f\|_\infty+V_a^b(f)\bigr)\|g-h\|_\infty.
\]
\end{twierdzenie}
\begin{proof}
Funkcję $g$ przybliżamy jednostajnie funkcjami schodkowymi $g_n$. Całki $I_n=(\KS)\int_a^b f\,\mathrm dg_n$ istnieją na mocy twierdzenia~\ref{tw:ks-step}. Z lematu~\ref{lem:ks-summation}, zastosowanego do $g_n-g_m$, otrzymujemy $|I_n-I_m|\le K_f\|g_n-g_m\|_\infty$, gdzie $K_f=2\|f\|_\infty+V_a^b(f)$. Istnieje zatem $I=\lim_nI_n$. Dla dowolnego podziału $\mathcal P$ mamy
$|S_g(f,\mathcal P)-I|\le K_f\|g-g_n\|_\infty+|S_{g_n}(f,\mathcal P)-I_n|+|I_n-I|$. Dla $\varepsilon>0$ wybieramy $n$ tak duże, aby $K_f\|g-g_n\|_\infty<\varepsilon/3$ i $|I_n-I|<\varepsilon/3$. Następnie wybieramy cechowanie dla całki $I_n$, dla którego środkowy składnik oszacowania jest mniejszy od $\varepsilon/3$. Wówczas każda drobna suma różni się od $I$ o mniej niż $\varepsilon$. To dowodzi istnienia całki. Z liniowości różnica całek jest całką względem $g-h$, które jest regulowane; stosując lemat \ref{lem:ks-summation} i przechodząc do granicy sum, otrzymujemy drugą tezę.
\end{proof}

\begin{uwaga}[Właściwa dziedzina obu twierdzeń]
Twierdzenia~\ref{tw:ks-reg-bv} i~\ref{tw:ks-bv-reg} są różne. Pierwsze pozwala na funkcję podcałkową regulowaną o nieskończonym wahaniu, o ile całkująca ma wahanie ograniczone; drugie pozwala na całkującą regulowaną o nieskończonym wahaniu, gdy funkcja podcałkowa ma wahanie ograniczone. Żadne z tych twierdzeń nie wymaga ciągłości obu funkcji ani braku wspólnych skoków.
\end{uwaga}

\begin{twierdzenie}[Konstruktywny rozkład funkcji o wahaniu ograniczonym]
\label{tw:ks-bv-split}
Dla $g\in\BV[a,b]$ zdefiniujmy
\[
 g_{\mathrm{sk}}(t)=
 \sum_{c\in(a,b]}\Delta^-g(c)\ind{[c,b]}(t)
 +\sum_{c\in[a,b)}\Delta^+g(c)\ind{(c,b]}(t).
\]
Oba szeregi są bezwzględnie i jednostajnie zbieżne, a funkcja $g_{\mathrm c}:=g-g(a)-g_{\mathrm{sk}}$ jest ciągła i ma wahanie ograniczone. Dla każdej regulowanej $f:[a,b]\to\R$ zachodzi
\[
 \begin{aligned}
 (\KS)\int_a^b f\,\mathrm dg
 ={}&(\KS)\int_a^b f\,\mathrm dg_{\mathrm c}
 +f(a)\Delta^+g(a)\\
 &+\sum_{c\in(a,b)}f(c)\bigl(\Delta^-g(c)+\Delta^+g(c)\bigr)
 +f(b)\Delta^-g(b).
 \end{aligned}
\]
W szczególności dla funkcji całkującej złożonej wyłącznie ze skoków całka jest szeregiem wartości $f$ w punktach skoków, niezależnie od ich liczby.
\end{twierdzenie}
\begin{proof}
Z lematu \ref{it:lem:ks-bv-regulated} funkcja $g$ jest
regulowana. Dla $a=b$ oba szeregi są puste, $g_{\mathrm c}=0$
i teza wynika z konwencji całki zerowej. Dalej $a<b$. Dla każdego skończonego zbioru punktów $E\subset[a,b]$ suma modułów wszystkich skoków jednostronnych $g$ w tych punktach nie przekracza $V_a^b(g)$. Aby to zobaczyć, wybieramy parami rozłączne, odpowiednio małe otoczenia punktów $E$, a następnie podziały uwzględniające punkty $c$ i po jednym punkcie z każdej dozwolonej strony $c$. W definicji wahania przechodzimy do granic jednostronnych. Biorąc supremum po skończonych $E$, otrzymujemy
$\sum_{c\in(a,b]}|\Delta^-g(c)|+\sum_{c\in[a,b)}|\Delta^+g(c)|\le V_a^b(g)$. Dla każdego $n\ge1$ zbiór punktów, w których
co najmniej jeden z dopuszczalnych skoków jednostronnych
ma moduł nie mniejszy niż $1/n$, jest skończony:
dla dowolnego skończonego podzbioru $E$ tego zbioru
mamy $\#E/n\le V_a^b(g)$. Zatem niezerowych skoków jest co najwyżej przeliczalnie wiele. Z kryterium Weierstrassa szeregi definiujące $g_{\mathrm{sk}}$ są bezwzględnie i jednostajnie zbieżne. Każdy skończony ogon ma wahanie nie większe od sumy modułów jego współczynników; przejście do granicy w sumie wariacyjnej dla dowolnego ustalonego podziału zachowuje tę nierówność. W szczególności $g_{\mathrm{sk}}\in\BV[a,b]$.

Także $g_{\mathrm{sk}}$ jest regulowana jako jednostajna granica funkcji schodkowych. Dla skończonych sum skok lewostronny w $c$ wynika
wyłącznie ze składnika $\Delta^-g(c)\chi_{[c,b]}$,
a prawostronny wyłącznie ze składnika
$\Delta^+g(c)\chi_{(c,b]}$, o ile dany składnik występuje.
Dla dwóch funkcji regulowanych $p,q$ mamy
$|p(c\pm)-q(c\pm)|\le\|p-q\|_\infty$ przez przejście
do granicy jednostronnej, a także
$|p(c)-q(c)|\le\|p-q\|_\infty$. Zatem
$|\Delta^\pm p(c)-\Delta^\pm q(c)|\le2\|p-q\|_\infty$
dla każdej dopuszczalnej strony. To uzasadnia przejście
do granicy w skokach jednostronnych. Otrzymujemy odpowiednio $\Delta^-g(c)$ i $\Delta^+g(c)$. Ponieważ $g$ jest regulowana, różnica $g_{\mathrm c}$ nie ma żadnych skoków jednostronnych i dlatego jest ciągła. Ma również wahanie ograniczone, jako różnica funkcji o wahaniu ograniczonym. Dla skończonych sum wzór na całkę wynika z twierdzenia~\ref{tw:ks-step}. Przy zwiększaniu liczby składników ogon ma wahanie ograniczone przez sumę modułów pominiętych skoków, a więc jego całka ma moduł nie większy niż $\|f\|_\infty$ razy ten ogon, na mocy lematu~\ref{lem:ks-var}. Przejście do granicy dowodzi wzoru.
\end{proof}

\subsection*{Wzory przez miarę Stieltjesa i rozkład Lebesgue'a}
\begin{twierdzenie}[Miara kanonicznie związana z funkcją $\BV$]
\label{tw:ks-measure}
Niech $g\in\BV[a,b]$. Dla $a=b$ przyjmujemy $\mu_g=0$; dalej zakładamy $a<b$. Definiujemy $G:[a,b]\to\R$ przez $G(t)=g(t+)-g(a)$ dla $a\le t<b$ oraz $G(b)=g(b)-g(a)$. Istnieje dokładnie jedna skończona \emph{miara ze znakiem borelowska} $\mu_g$ na $[a,b]$ spełniająca $\mu_g([a,t])=G(t)$ dla każdego $t\in[a,b]$. Jej atomy są dane wzorami
\[
 \mu_g(\{a\})=\Delta^+g(a),\qquad
 \mu_g(\{c\})=\Delta g(c)\ \ (a<c<b),\qquad
 \mu_g(\{b\})=\Delta^-g(b).
\]
Jeśli $f$ jest regulowana, to $f$ jest całkowalna względem $|\mu_g|$ i
\[
 (\KS)\int_a^b f\,\mathrm dg=\int_{[a,b]}f(t)\,\mathrm d\mu_g(t),
 \qquad |\mu_g|([a,b])\le V_a^b(g).
\]
\end{twierdzenie}
\begin{proof}
Niech $V(t)=V_a^t(g)$ oraz $U(t)=(V(t)+g(t)-g(a))/2$, $W(t)=(V(t)-g(t)+g(a))/2$. Z addytywności wahania i nierówności $V(v)-V(u)\ge|g(v)-g(u)|$ wynika, że $U,W$ są niemalejące, zerują się w $a$ oraz $U-W=g-g(a)$. Mają granice jednostronne. Dla $U$ definiujemy na prostej niemalejącą funkcję prawostronnie ciągłą $A_U$, równą $0$ dla $t<a$, $U(t+)$ dla $a\le t<b$, a $U(b)$ dla $t\ge b$. Dla $a\le t<b$ zachodzi $U(t+)=\inf_{t<s\le b}U(s)$; z tego wzoru i monotoniczności wynika $\lim_{u\downarrow t}U(u+)=U(t+)$. W $b$ prawostronna ciągłość wynika ze stałego przedłużenia. Analogicznie definiujemy $A_W$. Na mocy \atlasref{ME-04::it:thm:stieltjes-construction}{ME-04: konstrukcja miary Lebesgue'a--Stieltjesa} ich miary $\nu_U,\nu_W$ są skupione na $[a,b]$, mają dystrybuanty $A_U,A_W$ i masy odpowiednio $U(b),W(b)$. Zatem $\mu_g=\nu_U-\nu_W$ spełnia $\mu_g([a,t])=G(t)$, również dla $t=a,b$. Jednoznaczność skończonej miary ze znakiem wynika z równości jej wartości na przedziałach $[a,t]$, które tworzą układ $\pi$ generujący zbiory borelowskie w $[a,b]$. Różnica dwóch takich miar jest zerowa na całym przedziale, a jej zbiory o masie zero tworzą układ Dynkina, więc różnica jest zerowa na wszystkich zbiorach borelowskich. Dla $a<c\le b$ zachodzi $G(c-)=g(c-)-g(a)$: jeśli $t<c$ jest dostatecznie bliski $c$, wszystkie wartości $g(s)$ dla $t<s<c$ są dowolnie bliskie $g(c-)$, a więc także ich granica $g(t+)$. Z ciągłości skończonej miary ze znakiem z dołu otrzymujemy $\mu_g(\{c\})=G(c)-G(c-)$; dla $c=a$ mamy $\mu_g(\{a\})=G(a)$. Podstawienie definicji $G$ daje wszystkie trzy wzory na atomy. Ponieważ $|\nu_U-\nu_W|\le\nu_U+\nu_W$, otrzymujemy
$|\mu_g|([a,b])\le U(b)+W(b)=V_a^b(g)$.

Dla funkcji $\ind{\{c\}}$ równość całek wynika
z lematu \ref{lem:ks-point} i wzorów na atomy.
Dla $a\le u<v\le b$ z ciągłości miary z dołu mamy
$\mu_g([a,v))=G(v-)$, stąd
\[
 \mu_g((u,v))=G(v-)-G(u)=g(v-)-g(u+).
\]
Dodając masy odpowiednich końców $u,v$, otrzymujemy
pozostałe trzy wzory z lematu \ref{lem:ks-indicators},
z uwzględnieniem konwencji na końcach $a,b$.
Zatem równość zachodzi dla wszystkich tych funkcji
charakterystycznych. Z liniowości wynika równość dla każdej funkcji schodkowej. Wybieramy funkcje schodkowe $f_n\to f$ jednostajnie. Z twierdzenia~\ref{tw:ks-reg-bv} i lematu~\ref{lem:ks-var} całki KS z $f_n$ zbiegają do całki KS z $f$, zaś całki względem $\mu_g$ zbiegają dzięki nierówności $|\int(f_n-f)\,\mathrm d\mu_g|\le\|f_n-f\|_\infty|\mu_g|([a,b])$. Stąd teza.
\end{proof}

\begin{uwaga}[Konwencja na końcach i równość prawie wszędzie]
Jeśli $g$ jest niemalejąca i prawostronnie ciągła na $[a,b]$, to $\mu_g(\{a\})=0$ i $\mu_g((s,t])=g(t)-g(s)$; jest to konwencja używana wcześniej przy całce Lebesgue'a--Stieltjesa. Dla funkcji, która nie jest prawostronnie ciągła, nie wolno stosować tego samego wzoru bez uwzględnienia przesunięcia do $g(t+)$. Równość $f=h$ prawie wszędzie względem miary Lebesgue'a nie wystarcza do równości całek KS względem $g$ z atomami: trzeba uwzględnić wartości w punktach, którym $\mu_g$ przypisuje masę.
\end{uwaga}

\begin{twierdzenie}[Całkowalność względem miary integratora niemalejącego]\label{ma1:kh:thm:ks-positive-measure-l1}
Niech $a<b$, $g:[a,b]\to\R$ będzie niemalejąca (bez założenia prawostronnej ciągłości), a $\mu_g$ będzie miarą z twierdzenia \ref{tw:ks-measure}, uzupełnioną na zbiorach miary zero. Jeśli $f:[a,b]\to\R$ ma skończone wartości, jest mierzalna względem tej miary i $\int_{[a,b]}|f|\,\mathrm d\mu_g<\infty$, to $f$ jest KS-całkowalna względem $g$ i
\[
 (\KS)\int_a^b f\,\mathrm dg=\int_{[a,b]} f\,\mathrm d\mu_g.
\]
\end{twierdzenie}
\begin{proof}
Dla podziału $P=\{a=x_0<\dots<x_n=b\}$ definiujemy dodatnie miary
\[
 \nu_i=\mu_g|_{(x_{i-1},x_i)}
 +\Delta^+g(x_{i-1})\delta_{x_{i-1}}
 +\Delta^-g(x_i)\delta_{x_i},
\]
gdzie $\delta_c$ oznacza miarę Diraca w $c$. Z definicji dystrybuanty $G$ w twierdzeniu \ref{tw:ks-measure}, dla $u<v$ mamy
$\mu_g((u,v))=g(v-)-g(u+)$. Stąd $\nu_i([a,b])=g(x_i)-g(x_{i-1})$.
Ponadto $\sum_i\nu_i=\mu_g$: wnętrza komórek są rozłączne, w każdym węźle wewnętrznym suma obu jednostronnych skoków daje pełną masę atomu, a w końcach przedziału pojawia się tylko odpowiedni skok jednostronny. W szczególności $\nu_i\le\mu_g$. Każda $\nu_i$ jest skupiona na $[x_{i-1},x_i]$. Dla dowolnego podziału znaczonego
\[
 S_g(f,\mathcal P)-\int f\,\mathrm d\mu_g
 =\sum_i\int (f(t_i)-f(x))\,\mathrm d\nu_i(x).
\]

Niech $\varepsilon>0$ oraz $T=\mu_g([a,b])$. Wybieramy $h>0$ z $hT<\varepsilon/3$. Dzielimy dziedzinę na zbiory mierzalne $E_k=\{kh\le f<(k+1)h\}$, $k\in\Z$. Niech $w_k=2^{-|k|}/3$; $\sum_k w_k=1$. Z bezwzględnej ciągłości całki $|f|$ względem
$\mu_g$ wybieramy $\eta>0$ tak, aby $\mu_g(A)<\eta$ implikowało $\int_A|f|\,\mathrm d\mu_g<\varepsilon/3$. Uzasadnijmy zastosowanie regularności do uzupełnionej miary. Dla każdego mierzalnego w uzupełnieniu $E_k$ istnieją zbiory borelowskie $B_k\subset E_k\subset H_k$ z $\mu_g(H_k\setminus B_k)=0$, z definicji uzupełnienia. W szczególności $\mu_g(H_k)=\mu_g(E_k)$. Z regularności zewnętrznej skończonej miary borelowskiej
na przestrzeni metrycznej $[a,b]$
(\atlasref{ME-20::thm:finite-metric-regular-20}{ME-20: regularność skończonych miar borelowskich})
wybieramy względnie otwarte $O_k\supset H_k$, dla których $\mu_g(O_k\setminus E_k)=\mu_g(O_k)-\mu_g(H_k)$. Możemy więc wybrać $O_k$ tak, aby
\[
 \mu_g(O_k\setminus E_k)<w_k
 \min\{\eta/2,\varepsilon/(6(1+|kh|+h))\}.
\]
Dla $t\in E_k$ dobieramy cechowanie tak, by jego otoczenie względne leżało w $O_k$. Dla komórki z takim znacznikiem błąd na $E_k$ nie przekracza $h$, a poza $E_k$ nie przekracza $|kh|+h+|f(x)|$. Oznaczmy $U=\bigcup_k(O_k\setminus E_k)$. Dzięki $\sum_i\nu_i=\mu_g$, dla ustalonego $k$ suma miar złych części komórek ze znacznikami w $E_k$ nie przekracza $\mu_g(O_k\setminus E_k)$, zaś suma całek $|f|$ na wszystkich złych częściach nie przekracza $\int_U|f|\,\mathrm d\mu_g$. Zatem na każdym drobnym podziale
\[
 |S_g(f,\mathcal P)-\int f\,\mathrm d\mu_g|
 \le hT+\sum_k(|kh|+h)\mu_g(O_k\setminus E_k)
     +\int_U|f|\,\mathrm d\mu_g<\varepsilon.
\]
Istotnie, środkowy składnik nie przekracza $\varepsilon/6$, a $\mu_g(U)\le\eta/2$. To dowodzi tezy, również gdy $T=0$.
\end{proof}
\begin{wniosek}[Dostateczny warunek miarowy dla integratora o ograniczonym wahaniu]\label{ma1:kh:cor:ks-variation-measure-l1}
Niech $g\in\BV[a,b]$, a $V(t)=V_a^t(g)$. Jeżeli $f$ ma skończone wartości i jest całkowalna względem uzupełnionej dodatniej miary $\mu_V$, to jest KS-całkowalna względem $g$ i
\[
 (\KS)\int_a^b f\,\mathrm dg=\int_{[a,b]}f\,\mathrm d\mu_g.
\]
\end{wniosek}
\begin{proof}
Dla $a=b$ miary $\mu_g$ i $\mu_V$ są zerowe, więc obie całki są zerowe. Dalej zakładamy $a<b$. Z rozkładu $U=(V+g-g(a))/2$, $W=(V-g+g(a))/2$ mamy $U,W$ niemalejące, $\mu_U+\mu_W=\mu_V$ oraz $\mu_U-\mu_W=\mu_g$ (równości na przedziałach $[a,t]$ i jednoznaczność miary). Zatem $\mu_U,\mu_W\le\mu_V$ i $|\mu_g|\le\mu_V$ na zbiorach borelowskich. Wyjaśnijmy przejście do uzupełnień: dla każdej z tych trzech miar $\nu$ z dominacji wynika, że borelowski zbiór $\mu_V$-zerowy jest również $\nu$-zerowy. Każdy zbiór mierzalny względem uzupełnionej $\mu_V$ różni się od zbioru borelowskiego o podzbiór takiego zbioru zerowego, a więc jest mierzalny także względem uzupełnionej $\nu$. Dominacja zachowuje się na tych zbiorach; przez przybliżenie nieujemnymi funkcjami prostymi otrzymujemy $\int|f|\,\mathrm d\nu\le\int|f|\,\mathrm d\mu_V<\infty$. W szczególności $f$ spełnia wszystkie założenia całkowalności względem $\mu_U$ i $\mu_W$. Stosujemy twierdzenie \ref{ma1:kh:thm:ks-positive-measure-l1} osobno do $U,W$, a następnie liniowość KS. Nie twierdzimy, że $\mu_V=|\mu_g|$: przeciwne skoki jednostronne w tym samym punkcie mogą się znosić w $\mu_g$. Przykład \ref{ma1:kh:ex:ks-cancelled-spikes} pokazuje, że bez dodatkowego warunku nie można zastąpić $\mu_V$ miarą $|\mu_g|$.
\end{proof}
\begin{lemat}[Integratory skupione w pojedynczych punktach]\label{ma1:kh:lem:ks-point-spikes}
Niech $r(a)=r(b)=0$, a $r(t)=0$ poza co najwyżej przeliczalnym zbiorem $C\subset(a,b)$. Załóżmy $\sum_{c\in C}|r(c)|<\infty$. Jeżeli $f:[a,b]\to\R$ ma skończone wartości oraz
\[
 \sum_{c\in C}|f(c)r(c)|<\infty,
\]
to $(\KS)\int_a^b f\,\mathrm dr=0$.
\end{lemat}
\begin{proof}
Ustalmy $\varepsilon>0$. Wybieramy skończony $C_0\subset C$ tak, aby $\sum_{c\in C\setminus C_0}|f(c)r(c)|<\varepsilon/4$. Dla $k\ge1$ wybieramy skończone $C_k\supset C_0$ zawarte w $C$, dla których
$\sum_{c\in C\setminus C_k}|r(c)|<\varepsilon 2^{-k-3}/k$.
Jeśli $k-1\le |f(t)|<k$, dobieramy $\delta(t)>0$ tak, aby otoczenie znacznika $t$ nie zawierało żadnego punktu $C_k\setminus\{t\}$. Jest to możliwe dzięki skończoności $C_k$.
Dla drobnego podziału z węzłami $x_i$ i znacznikami $t_i$ dokładne sumowanie daje
\[
 S_r(f,\mathcal P)=\sum_{i=1}^{n-1}r(x_i)(f(t_i)-f(t_{i+1})).
\]
Wstawiamy i odejmujemy w każdym nawiasie $f(x_i)$. Składniki z $t_j=x_i$ znikają; składnik z $t_j\ne x_i$ i $r(x_i)\ne0$ jest możliwy jedynie dla $x_i\in C\setminus C_k$, gdzie $k-1\le |f(t_j)|<k$. Każdy węzeł pojawia się co najwyżej dwukrotnie. Wkład pochodzący od wartości $f$ w znacznikach ograniczamy więc przez
$2\sum_{k\ge1}k\sum_{c\in C\setminus C_k}|r(c)|\le\varepsilon/4$,
a wkład wartości w węzłach przez
$2\sum_{c\in C\setminus C_0}|f(c)r(c)|<\varepsilon/2$.
W rezultacie $|S_r(f,\mathcal P)|<\varepsilon$, co dowodzi tezy. Dowód nie wymaga mierzalności $f$ względem miary Lebesgue'a.
\end{proof}
\begin{twierdzenie}[Kryterium miarowe z kontrolą punktowych odchyleń]\label{ma1:kh:thm:ks-signed-measure-l1}
Niech $g\in\BV[a,b]$ i niech $f$ ma skończone wartości oraz będzie całkowalna względem uzupełnionej miary $|\mu_g|$. Jeżeli dodatkowo
\[
 \sum_{c\in(a,b)}|f(c)|\,|g(c)-g(c+)|<\infty,
\]
to $f$ jest KS-całkowalna względem $g$ i jej całka jest równa $\int_{[a,b]}f\,\mathrm d\mu_g$. W szczególności dla prawostronnie ciągłego integratora warunek dodatkowy znika.
\end{twierdzenie}
\begin{proof}
Dla $a=b$ obie całki są zerowe, więc dalej zakładamy $a<b$. Definiujemy $h(a)=g(a)$, $h(b)=g(b)$ oraz $h(t)=g(t+)$ dla $a<t<b$. Mamy $r=g-h$ zerową na końcach, a jej niezerowe wartości to $-\Delta^+g(c)$, których moduły tworzą szereg zbieżny na mocy twierdzenia \ref{tw:ks-bv-split}. Ponadto $\mu_h=\mu_g$: obie funkcje mają te same granice prawostronne i te same wartości w końcach. Dokładniej, niech $\mu_g=\mu^+-\mu^-$ będzie rozkładem Jordana miary. Definiujemy niemalejące funkcje $h_\pm$ przez $h_\pm(a)=0$ oraz $h_\pm(t)=\mu^\pm([a,t])$ dla $a<t\le b$. Wówczas $h=g(a)+h_+-h_-$ (w szczególności $h\in\BV$), a z definicji miary integratora $\mu_{h_\pm}=\mu^\pm$, także przy ewentualnym atomie w $a$. Całkowalność względem $|\mu_g|=\mu^++\mu^-$ daje całkowalność względem każdej z tych miar. Twierdzenie \ref{ma1:kh:thm:ks-positive-measure-l1} zastosowane do $h_\pm$ i liniowość dają $(\KS)\int f\,\mathrm dh=\int f\,\mathrm d\mu_g$. Lemat \ref{ma1:kh:lem:ks-point-spikes} daje $(\KS)\int f\,\mathrm dr=0$, co kończy dowód.
\end{proof}
\begin{przyklad}[Sama całkowalność względem $|\mu_g|$ nie wystarcza]\label{ma1:kh:ex:ks-cancelled-spikes}
Na $[0,1]$ połóżmy $c_n=2^{-n}$, $g(c_n)=2^{-n}$, $f(c_n)=2^n$, a poza tymi punktami $g=f=0$. Funkcja $g$ ma ograniczone wahanie: jest sumą funkcji $2^{-n}\chi_{\{c_n\}}$, których wahania sumują się do $2\sum_n2^{-n}=2$. Granice prawostronne i lewostronne $g$ wynoszą zero, również w punkcie skupienia $0$, a $g(0)=g(1)=0$. Zatem $\mu_g=0$, więc $f$ jest całkowalna względem $|\mu_g|$ z całką zero.
Przypuśćmy, że istnieje całka KS z $f$ względem $g$ i niech $F$ będzie jej pierwotną. Z twierdzenia \ref{tw:ks-primitive} wynika istnienie skończonej granicy prawostronnej $F(0+)$, a dla każdego $n$ zachodzi
\[
 F(c_n)-F(c_n-)=f(c_n)\Delta^-g(c_n)=1.
\]
Jest to niemożliwe: jeśli $F(0+)=L$, to dla dostatecznie dużych $n$ zarówno $F(c_n)$, jak i $F(c_n-)$ należą do $[L-1/4,L+1/4]$ (drugą wartość otrzymujemy, przechodząc do granicy w punktach dodatnich mniejszych od $c_n$). Ich różnica ma moduł co najwyżej $1/2$. Sprzeczność dowodzi nieistnienia całki KS.
\end{przyklad}
\begin{twierdzenie}[Porównanie dwóch kryteriów miarowych]\label{ma1:kh:thm:ks-measure-criteria-equivalence}
Niech $g\in\BV[a,b]$, $V(t)=V_a^t(g)$ i niech $f$ ma skończone wartości. Całkowalność $f$ względem uzupełnionej miary $\mu_V$ jest równoważna łącznemu spełnieniu warunków
\[
 f\in\mathcal L^1(|\mu_g|),\qquad
 \sum_{c\in(a,b)}|f(c)|\,|\Delta^+g(c)|<\infty,
\]
gdzie $\mathcal L^1$ oznacza klasę konkretnych funkcji całkowalnych względem uzupełnionej miary, bez przechodzenia do klas równoważności.
\end{twierdzenie}
\begin{proof}
Dla funkcji $p,q\in\BV[a,b]$ oznaczmy przez $V_p,V_q$ ich funkcje wahania od $a$. Z nierówności trójkąta dla sum wariacyjnych i addytywności wahania wynika, że $V_p+V_q-V_{p+q}$ jest niemalejąca i zeruje się w $a$. Liniowość przyporządkowania integratorowi jego miary wynika z jednoznaczności w twierdzeniu \ref{tw:ks-measure}; miara integratora niemalejącego jest dodatnia. Zatem $\mu_{V_{p+q}}\le\mu_{V_p}+\mu_{V_q}$.

Niech $h,r$ będą funkcjami z dowodu twierdzenia \ref{ma1:kh:thm:ks-signed-measure-l1}, tak że $g=h+r$. Rozkład $h=g(a)+h_+-h_-$ z tego dowodu daje $\mu_{V_h}\le\mu_{h_+}+\mu_{h_-}=|\mu_g|$, ponieważ dla niemalejącego integratora jego funkcja wahania różni się od niego o stałą, a $V_{-h_-}=V_{h_-}$. Dla $r$ zachodzi
\[
 \mu_{V_r}\le 2\sum_{c\in(a,b)}|r(c)|\,\delta_c.
\]
Istotnie, dla $u<v$ suma wariacyjna $r$ na $[u,v]$ nie przekracza $|r(u)|+2\sum_{u<c<v}|r(c)|+|r(v)|$, gdyż każdy wewnętrzny węzeł występuje w co najwyżej dwóch składnikach. Funkcja
\[
 A(t)=\sum_{c\in(a,b)}|r(c)|\bigl(\chi_{[c,b]}(t)+\chi_{(c,b]}(t)\bigr)
\]
spełnia zatem $A(v)-A(u)\ge V_r(v)-V_r(u)$, czyli $A-V_r$ jest niemalejąca. Miara $A$ jest równa $2\sum_c|r(c)|\delta_c$: wynika to z jej dystrybuanty i zbieżności jednostajnej szeregu. Otrzymujemy
\[
 |\mu_g|\le\mu_V\le|\mu_g|+2\sum_{c\in(a,b)}|\Delta^+g(c)|\,\delta_c.
\]
Pierwsza nierówność była wykazana we wniosku \ref{ma1:kh:cor:ks-variation-measure-l1}. Ponadto dla każdego $c\in(a,b)$ mamy $\mu_V(\{c\})=V(c+)-V(c-)\ge V(c+)-V(c)\ge|\Delta^+g(c)|$. Pierwsza nierówność wynika z monotoniczności $V$, a druga z $V(v)-V(c)\ge|g(v)-g(c)|$ dla $v>c$ po przejściu do granicy przy $v\downarrow c$.

Jeśli $f$ jest całkowalna względem $\mu_V$, pierwsza nierówność daje całkowalność względem $|\mu_g|$, a oszacowanie atomów daje zbieżność szeregu. Odwrotnie, jeśli spełnione są dwa warunki z tezy, funkcja $f$ jest mierzalna względem uzupełnienia miary z prawej strony ostatniej podwójnej nierówności: poza przeliczalnym zbiorem jej dodatkowych atomów zgadza się prawie wszędzie względem $|\mu_g|$ z funkcją borelowską, a w tych atomach możemy nadać funkcji borelowskiej wartości $f$. Pozostały zbiór wyjątkowy jest zerowy dla obu składników miary. Z dominacji wynika mierzalność względem uzupełnionej $\mu_V$ i skończoność całki modułu. To kończy dowód.
\end{proof}
\begin{twierdzenie}[Dokładna różnica między miarą wahania i wahaniem miary]\label{ma1:kh:thm:ks-variation-measure-defect}
Niech $a<b$, $g\in\BV[a,b]$ i $V(t)=V_a^t(g)$. Zachodzi równość skończonych dodatnich miar
\[
 \mu_V=|\mu_g|+\sum_{c\in(a,b)}d_c\,\delta_c,
 \qquad
 d_c=|\Delta^-g(c)|+|\Delta^+g(c)|
      -|\Delta^-g(c)+\Delta^+g(c)|.
\]
Współczynniki $d_c$ są nieujemne i sumowalne. W szczególności $\mu_V=|\mu_g|$ wtedy i tylko wtedy, gdy w każdym punkcie wewnętrznym oba skoki jednostronne mają zgodne znaki lub przynajmniej jeden z nich jest zerowy. Równoważnie $g(c)$ leży pomiędzy $g(c-)$ i $g(c+)$ dla każdego $c\in(a,b)$. Dla funkcji o skończonych wartościach dokładnym warunkiem przynależności do $\mathcal L^1(\mu_V)$ jest przynależność do $\mathcal L^1(|\mu_g|)$ oraz $\sum_c d_c|f(c)|<\infty$; miary rozumiemy jako uzupełnione.
\end{twierdzenie}
\begin{proof}
Najpierw obliczymy skoki funkcji wahania. Ustalmy $\varepsilon>0$. Dla $c<b$ wybieramy podział $c=y_0<y_1<\cdots<y_m=b$, którego suma wariacyjna jest większa niż $V_c^b(g)-\varepsilon$. Jeśli $c<t<y_1$, to zastępując w nim początek $c$ przez $t$, zmniejszamy sumę o co najwyżej $|g(t)-g(c)|$ na mocy odwrotnej nierówności trójkąta. Otrzymujemy
\[
 V_t^b(g)>V_c^b(g)-\varepsilon-|g(t)-g(c)|.
\]
Z addytywności wahania (lemat \ref{ma1:kh:lem:variation-additivity}) i podziału złożonego z samych końców wynika
\[
 |g(t)-g(c)|\le V(t)-V(c)
 < |g(t)-g(c)|+\varepsilon .
\]
Przechodząc z $t\downarrow c$ i następnie $\varepsilon\downarrow0$, dostajemy $\Delta^+V(c)=|\Delta^+g(c)|$. Analogicznie dla $c>a$ wybieramy podział $[a,c]$ z sumą większą niż $V_a^c(g)-\varepsilon$. Dla $t<c$ leżącego za przedostatnim węzłem zastępujemy jego ostatni węzeł $c$ przez $t$, uzyskując $V(t)>V(c)-\varepsilon-|g(c)-g(t)|$. Ponieważ $V(c)-V(t)\ge |g(c)-g(t)|$, przejście z $t\uparrow c$ daje $\Delta^-V(c)=|\Delta^-g(c)|$. Dowody obejmują też odpowiednie skoki w końcach przedziału.

Z dowodu twierdzenia \ref{ma1:kh:thm:ks-measure-criteria-equivalence} mamy nierówności miar
\[
 0\le\mu_V-|\mu_g|
 \le 2\sum_{c\in(a,b)}|\Delta^+g(c)|\,\delta_c .
\]
Zatem dodatnia miara po lewej jest skupiona na co najwyżej przeliczalnym zbiorze punktów wewnętrznych. W każdym z nich jej masa wynosi
\[
 \mu_V(\{c\})-|\mu_g|(\{c\})
 =|\Delta^-g(c)|+|\Delta^+g(c)|
  -|\Delta^-g(c)+\Delta^+g(c)|,
\]
na mocy wzorów na atomy z twierdzenia \ref{tw:ks-measure} i obliczonych skoków $V$. Dla miary ze znakiem jej wahanie na zbiorze jednopunktowym jest modułem masy tego zbioru, co wynika bezpośrednio z definicji wahania miary. Zbiór $C=\{c\in(a,b):\Delta^+g(c)\ne0\}$ jest co najwyżej przeliczalny na mocy twierdzenia \ref{tw:ks-bv-split}. Dla dowolnego borelowskiego $B\subseteq[a,b]$ miara różnicowa ma wartość równą sumie swoich mas na $B\cap C$, gdyż jest zerowa na $B\setminus C$. W szczególności nie ma dodatkowego składnika na końcach przedziału. Przeliczalna addytywność daje żądaną równość. Sumowalność $d_c$ wynika ze skończoności $\mu_V$; ponadto dla liczb rzeczywistych $p,q$ zachodzi $|p|+|q|-|p+q|=0$ dokładnie wtedy, gdy $pq\ge0$. Stąd kryterium równości miar, a po podstawieniu $p=g(c)-g(c-)$, $q=g(c+)-g(c)$ także równoważny warunek na wartość $g(c)$.

Jeżeli $f$ jest całkowalna względem $\mu_V$, to dominacja $|\mu_g|\le\mu_V$ i powyższa równość miar dają oba warunki. Odwrotnie, z mierzalności $f$ względem uzupełnienia $|\mu_g|$ wybieramy borelowskiego reprezentanta $f$ poza zbiorem $|\mu_g|$-zerowym i korygujemy go w przeliczalnym zbiorze $\{c:d_c>0\}$, nadając tam rzeczywiste wartości $f$. Skorygowany reprezentant pozostaje borelowski, a zbiór, na którym różni się od $f$, jest zerowy dla obu składników rozkładu $\mu_V$. Dowodzi to mierzalności względem uzupełnionej $\mu_V$. Z addytywności całek nieujemnych względem sumy miar otrzymujemy
\[
 \int |f|\,\mathrm d\mu_V
 =\int |f|\,\mathrm d|\mu_g|+\sum_c d_c|f(c)|<\infty .
\]
\end{proof}
\begin{przyklad}[Nadmiar wahania w pojedynczym punkcie]\label{ma1:kh:ex:ks-variation-defect}
Na $[0,1]$ przyjmijmy $c=1/2$, $g(x)=0$ dla $x<c$, $g(c)=3$ i $g(x)=1$ dla $x>c$. Wówczas $\Delta^-g(c)=3$, $\Delta^+g(c)=-2$ oraz $\mu_g=\delta_c$. Funkcja wahania $V$ przyjmuje kolejno wartości $0,3,5$, więc $\mu_V=5\delta_c$. Różnica $\mu_V-|\mu_g|=4\delta_c$ jest dokładnie opisana przez twierdzenie \ref{ma1:kh:thm:ks-variation-measure-defect}. Dla każdej funkcji o skończonych wartościach całka KS wynosi $f(c)$ (wzór na skoki z twierdzenia \ref{tw:ks-step}), choć w oszacowaniach opartych na wahaniu integratora występuje masa $5$ zamiast $1$.
\end{przyklad}
\begin{twierdzenie}[Punktowe odchylenia z jednym punktem skupienia]\label{ma1:kh:thm:ks-discrete-spikes}
Niech $a<c_{n+1}<c_n<b$, $c_n\downarrow a$, $r(c_n)=r_n$ oraz $r=0$ poza $\{c_n:n\ge1\}$. Załóżmy $r_n\to0$. Dla dowolnej funkcji $f$ o skończonych wartościach zachodzi
\[
 f\in\mathrm{KS}(r)\quad\Longleftrightarrow\quad
 f(c_n)r_n\longrightarrow0.
\]
Jeżeli całka istnieje, jest równa zeru. Nie zakładamy bezwzględnej sumowalności $r_n$ ani $f(c_n)r_n$.
\end{twierdzenie}
\begin{proof}
Konieczność wynika z twierdzenia \ref{tw:ks-primitive}: pierwotna $F$ ma granicę prawostronną w $a$, ponieważ $r(a+)=0$, a jej skok lewostronny w $c_n$ jest równy $f(c_n)r_n$. Istnienie granicy $F(a+)$ wymusza zbieżność tych skoków do zera, jak w przykładzie \ref{ma1:kh:ex:ks-cancelled-spikes}.

Dla dostateczności niech $\varepsilon>0$. Zbiór $D=\{a\}\cup\{c_n:n\ge1\}$ jest domknięty. Dla $t\notin D$ wybieramy $\delta(t)<\operatorname{dist}(t,D)$. Dla $t=c_n$ wybieramy $\delta(t)>0$ tak małe, aby otoczenie zawierało spośród $D$ jedynie $c_n$; to możliwe, gdyż $c_n$ jest punktem izolowanym w $D$. W punkcie $a$ wybieramy $\delta(a)>0$ tak małe, aby dla wszystkich $c_n<a+\delta(a)$ zachodziło
\[
 (|f(a)|+|f(c_n)|)|r_n|<\varepsilon.
\]
Istnieje takie cechowanie, ponieważ $r_n\to0$ i $f(c_n)r_n\to0$.

W drobnym podziale pierwsza komórka musi mieć znacznik $a$, a żadna inna nie może go mieć. Każda pozostała komórka zawierająca $c_n$ musi mieć znacznik $c_n$. Z tożsamości sumowania przez części (przy $r(a)=r(b)=0$)
\[
 S_r(f,\mathcal P)=\sum_{i=1}^{m-1}r(x_i)(f(t_i)-f(t_{i+1}))
\]
wszystkie składniki są zerowe poza ewentualnym składnikiem dla prawego końca pierwszej komórki, gdy ten koniec jest pewnym $c_n$. W takim przypadku jedyny składnik to $r_n(f(a)-f(c_n))$, którego moduł jest mniejszy od $\varepsilon$ z wyboru $\delta(a)$. Jeżeli prawy koniec pierwszej komórki nie należy do zbioru $\{c_n:n\ge1\}$, suma wynosi zero. W obu przypadkach definicja całki daje wartość zero.
\end{proof}
\begin{przyklad}[Bezwzględna sumowalność punktowych odchyleń nie jest konieczna]\label{ma1:kh:ex:ks-nonsummable-spikes}
Na $[0,1]$ niech $c_n=2^{-n}$, $g(c_n)=2^{-n}$, $f(c_n)=2^n/n$ i $f=g=0$ poza $\{c_n:n\ge1\}$. Jak w przykładzie \ref{ma1:kh:ex:ks-cancelled-spikes}, $g\in\BV[0,1]$ i $\mu_g=0$. Ponieważ $f(c_n)g(c_n)=1/n\to0$, twierdzenie \ref{ma1:kh:thm:ks-discrete-spikes} daje $(\KS)\int_0^1 f\,\mathrm dg=0$. Natomiast
\[
 \sum_{n\ge1}|f(c_n)|\,|\Delta^+g(c_n)|=\sum_{n\ge1}\frac1n=\infty.
\]
Z twierdzenia \ref{ma1:kh:thm:ks-measure-criteria-equivalence} wynika też brak całkowalności względem $\mu_V$. Obydwa równoważne kryteria miarowe są więc warunkami dostatecznymi, lecz nie koniecznymi, nawet dla integratora o ograniczonym wahaniu z zerową miarą ze znakiem.
\end{przyklad}
\begin{twierdzenie}[Kryterium bezwzględnej całkowalności KS dla integratora niemalejącego]\label{ma1:kh:thm:ks-absolute-monotone}
Niech $g:[a,b]\to\R$ będzie niemalejąca, a $f:[a,b]\to\R$ będzie mierzalna względem uzupełnionej miary $\mu_g$. Jeżeli $f\ge0$ prawie wszędzie względem $\mu_g$ i $f$ jest KS-całkowalna względem $g$, to $f\in\mathcal L^1(\mu_g)$. W konsekwencji dla takich mierzalnych funkcji zachodzi
\[
 f\in\mathcal L^1(\mu_g)\quad\Longleftrightarrow\quad
 f,|f|\in\mathrm{KS}(g).
\]
Nie zakładamy prawostronnej ciągłości $g$. Wartości funkcji są skończone, a w obu przypadkach całka KS jest równa całce względem $\mu_g$.
\end{twierdzenie}
\begin{proof}
Najpierw zauważmy, że jeśli $N$ jest zbiorem $\mu_g$-zerowym, to dowolna skończona funkcja zerowa poza $N$ należy do $\mathcal L^1(\mu_g)$ (miara jest uzupełniona) i ma całkę KS zero z twierdzenia \ref{ma1:kh:thm:ks-positive-measure-l1}. Można więc zmienić $f$ na zbiorze wyjątkowym tak, aby $f\ge0$ wszędzie, bez zmiany jej całkowalności KS ani wartości całki. Dla niemalejącego $g$ każda suma KS funkcji nieujemnej jest nieujemna, ponieważ wszystkie przyrosty $g$ są nieujemne. Wybierając drobny podział dla dowolnej dokładności, otrzymujemy dodatniość całki; z liniowości wynika monotoniczność na klasie funkcji KS-całkowalnych.
Dla $h_n=\min(f,n)$ mamy $h_n\in\mathcal L^1(\mu_g)$ ze skończoności miary oraz mierzalności $f$. Twierdzenie \ref{ma1:kh:thm:ks-positive-measure-l1} daje całkowalność KS funkcji $h_n$ i zgodność wartości, a monotoniczność daje
\[
 \int h_n\,\mathrm d\mu_g=(\KS)\int_a^b h_n\,\mathrm dg
 \le(\KS)\int_a^b f\,\mathrm dg<\infty.
\]
Z twierdzenia o zbieżności monotonicznej, ponieważ $h_n\uparrow f$, otrzymujemy $\int f\,\mathrm d\mu_g<\infty$. Przywrócenie wartości na zbiorze wyjątkowym nie zmienia tego wniosku. Zastosowanie wyniku do $|f|$ dowodzi implikacji od prawej do lewej w równoważności. Druga implikacja oraz zgodność wartości wynikają z twierdzenia \ref{ma1:kh:thm:ks-positive-measure-l1}.
\end{proof}
\begin{wniosek}[Twierdzenia graniczne KS z majorantami KS]\label{ma1:kh:cor:ks-order-limits-monotone}
Niech $g$ będzie niemalejąca, a $f_n:[a,b]\to\R$ będą mierzalne względem uzupełnionej $\mu_g$ i KS-całkowalne względem $g$. Niech $f_n\to f$ prawie wszędzie względem $\mu_g$, z granicą o skończonych wartościach. Wówczas $f$ jest KS-całkowalna i całki $f_n$ zbiegają do całki $f$, jeżeli zachodzi jeden z warunków:
\begin{enumerate}[label=(\alph*)]
\item $f_n\uparrow f$ prawie wszędzie i całki KS funkcji $f_n$ są ograniczone z góry;
\item $p\le f_n\le q$ prawie wszędzie dla dwóch funkcji $p,q$ mierzalnych względem $\mu_g$ i KS-całkowalnych względem $g$.
\end{enumerate}
W obu przypadkach różnice $f_n-f$ są całkowalne względem $\mu_g$ i $\int|f_n-f|\,\mathrm d\mu_g\to0$.
\end{wniosek}
\begin{proof}
Zbieramy wszystkie zbiory wyjątkowe w jeden zbiór $\mu_g$-zerowy. Miara jest uzupełniona, więc $f$ jest mierzalna i możemy nadać wszystkim funkcjom wartość zero na tym zbiorze bez zmiany istniejących całek KS, na mocy dowodu twierdzenia \ref{ma1:kh:thm:ks-absolute-monotone}. W (a) funkcje $u_n=f_n-f_1$ są nieujemne, mierzalne i KS-całkowalne. To samo twierdzenie daje $u_n\in\mathcal L^1(\mu_g)$. Ich całki są ograniczone z góry przez założenie i liniowość KS. Zbieżność monotoniczna daje $u=f-f_1\in\mathcal L^1(\mu_g)$ oraz $\int(u-u_n)\,\mathrm d\mu_g\to0$. W (b) funkcje $u_n=f_n-p$ i $q-p$ są nieujemne, mierzalne i KS-całkowalne, więc są całkowalne względem $\mu_g$. Stosujemy zbieżność zmajoryzowaną do $0\le u_n\le q-p$, otrzymując $u=f-p\in\mathcal L^1(\mu_g)$ i $\int|u_n-u|\,\mathrm d\mu_g\to0$. W obu przypadkach twierdzenie \ref{ma1:kh:thm:ks-positive-measure-l1} daje całkowalność KS funkcji $u$, a liniowość daje całkowalność $f$ i zbieżność całek. Różnica $f_n-f=u_n-u$ dowodzi ostatniej tezy. Przywracamy początkowe wartości wszystkich funkcji na wspólnym zbiorze zerowym: całkowalność KS i wartości całek pozostają niezmienione na mocy pierwszego kroku dowodu twierdzenia \ref{ma1:kh:thm:ks-absolute-monotone}, a całkowalność i normy różnic względem uzupełnionej miary $\mu_g$ nie zmieniają się. Nie wymagamy, aby $p,q$ czy same $f_n$ były całkowalne względem $\mu_g$.
\end{proof}
\begin{wniosek}[Twierdzenia graniczne w klasie miarowej KS]\label{ma1:kh:cor:ks-measure-limits}
Przy $g\in\BV[a,b]$ oraz $V(t)=V_a^t(g)$ niech $f_n,f$ mają skończone wartości, będą mierzalne względem uzupełnionej miary $\mu_V$, niech $f_n\to f$ prawie wszędzie względem $\mu_V$ i $|f_n|\le h$ prawie wszędzie względem $\mu_V$ dla pewnego $h\in\mathcal L^1(\mu_V)$. Wówczas wszystkie całki KS istnieją, a całki $f_n$ zbiegają do całki $f$. Dla niemalejącego $g$ analogiczna teza zachodzi, gdy $f_n$ są całkowalne względem $\mu_g$, rosną do $f$ prawie wszędzie względem $\mu_g$ i ich całki są ograniczone z góry.
\end{wniosek}
\begin{proof}
W pierwszym przypadku $|f|\le h$ prawie wszędzie, przez przejście do granicy poza przeliczalną sumą zbiorów wyjątkowych. Stąd $f_n,f\in\mathcal L^1(\mu_V)$ oraz $|f_n-f|\le2h$ prawie wszędzie. Stosujemy twierdzenie o zbieżności zmajoryzowanej względem $\mu_V$ do $|f_n-f|$, a następnie wniosek \ref{ma1:kh:cor:ks-variation-measure-l1} i nierówność
$|\int(f_n-f)\,\mathrm d\mu_g|\le\int |f_n-f|\,\mathrm d\mu_V$.
W drugim przypadku stosujemy zbieżność monotoniczną do nieujemnych $f_n-f_1$ względem $\mu_g$ (ich całki są ograniczone z góry). Otrzymujemy całkowalność $f$ i zbieżność całek, po czym używamy twierdzenia \ref{ma1:kh:thm:ks-positive-measure-l1}. Nie są to twierdzenia o dowolnych funkcjach KS-całkowalnych: zakładamy odpowiednią całkowalność miarową.
\end{proof}

\begin{twierdzenie}[Wzór: część absolutnie ciągła, osobliwa i skoki]
\label{tw:ks-decomp}
Niech $g\in\BV[a,b]$ i $f:[a,b]\to\R$ będzie regulowana. Niech rozkład Lebesgue'a miary ze znakiem $\mu_g$ względem $\lambda_1$ ma postać
$\mu_g=q\lambda_1+\mu_{\mathrm{sc}}+\mu_{\mathrm{at}}$, gdzie $q\in L^1([a,b],\lambda_1)$, $\mu_{\mathrm{sc}}$ jest bezatomową miarą ze znakiem osobliwą względem $\lambda_1$, a $\mu_{\mathrm{at}}$ jest częścią czysto atomową. Wtedy
\[
 \begin{aligned}
 (\KS)\int_a^b f\,\mathrm dg
 ={}&\int_a^b f(t)q(t)\,\mathrm d\lambda_1(t)
 +\int_{[a,b]}f\,\mathrm d\mu_{\mathrm{sc}}\\
 &+f(a)\Delta^+g(a)+\sum_{c\in(a,b)}f(c)\Delta g(c)
 +f(b)\Delta^-g(b).
 \end{aligned}
\]
Szereg jest bezwzględnie zbieżny. Jeżeli $g$ nie ma osobliwej części ciągłej, drugi składnik znika; jeśli $g$ jest absolutnie ciągła, znikają również wszystkie składniki skokowe i $q=g'$ prawie wszędzie.
\end{twierdzenie}
\begin{proof}
Z twierdzenia~\ref{tw:ks-measure} całka KS jest całką względem skończonej miary ze znakiem $\mu_g$. Stosujemy rozkład Lebesgue'a jej części bezatomowej względem $\lambda_1$ i liniowość całki. Każda skończona miara ze znakiem ma co najwyżej przeliczalnie wiele atomów niezerowej masy, a suma modułów ich mas nie przekracza $|\mu_g|([a,b])\le V_a^b(g)$. Stąd $\sum|f(c)\Delta g(c)|\le\|f\|_\infty V_a^b(g)$ i bezwzględna zbieżność. Ostatnia część wynika ze wzoru Newtona--Leibniza dla funkcji absolutnie ciągłej $g(t)=g(a)+\int_a^t g'(s)\,\mathrm d\lambda_1(s)$.
\end{proof}

\begin{wniosek}[Wzór dla funkcji absolutnie ciągłej z impulsami]
\label{wn:ks-ac-jumps}
Niech $h\in\AC[a,b]$, $a<c_1<\cdots<c_m<b$, $A_j,B_j\in\R$ oraz
$g(t)=h(t)+\sum_{j=1}^m\bigl(A_j\ind{[c_j,b]}(t)+B_j\ind{(c_j,b]}(t)\bigr)$. Dla każdej regulowanej $f$ zachodzi
\[
 (\KS)\int_a^b f\,\mathrm dg
 =\int_a^b f(t)h'(t)\,\mathrm d\lambda_1(t)
 +\sum_{j=1}^m(A_j+B_j)f(c_j).
\]
Gdy $h\in C^1[a,b]$, pierwsza całka jest zwykłą całką Riemanna z funkcji regulowanej $fh'$.
\end{wniosek}
\begin{proof}
Dla części $h$ stosujemy twierdzenie~\ref{tw:ks-decomp}; dla każdej funkcji skokowej stosujemy twierdzenie~\ref{tw:ks-step}. Liniowość względem funkcji całkującej kończy dowód.
\end{proof}

\begin{przyklad}[Skończona liczba impulsów: pełny rachunek]
Na $[0,1]$ niech $g(t)=t+3\ind{[1/3,1]}(t)-2\ind{(2/3,1]}(t)$ oraz $f(t)=t^2$. Ponieważ funkcja $t\mapsto t$ jest absolutnie ciągła, a skoki $g$ mają wielkości odpowiednio $3$ i $-2$, otrzymujemy
\[
 (\KS)\int_0^1 t^2\,\mathrm dg(t)
 =\int_0^1 t^2\,\mathrm d\lambda_1(t)
 +3\left(\frac13\right)^2-2\left(\frac23\right)^2
 =\frac13+\frac13-\frac89=-\frac29.
\]
Wartość $g(2/3)$ jest tutaj inna niż $g((2/3)+)$, ale wzór na całkę korzysta z całkowitego skoku $\Delta g(2/3)$.
\end{przyklad}

\begin{przyklad}[Dwa atomy na końcach przedziału]
Na $[0,1]$ niech $g(t)=\ind{(0,1]}(t)+2\ind{\{1\}}(t)$ i $f(t)=1+t$. Funkcja $g$ ma atom o masie $1$ w punkcie $0$ i atom o masie $2$ w punkcie $1$, mimo że wewnątrz przedziału jest stała. Z twierdzenia~\ref{tw:ks-step} o całkowaniu względem funkcji schodkowej (lub z twierdzenia~\ref{tw:ks-measure} i wzorów na masy atomów) otrzymujemy
\[
 (\KS)\int_0^1 f\,\mathrm dg=f(0)\cdot1+f(1)\cdot2=5.
\]
Ten rachunek pokazuje, dlaczego w twierdzeniach trzeba osobno określić sposób traktowania końców przedziału.
\end{przyklad}

\begin{przyklad}[Nieskończenie wiele impulsów i oszacowanie ogona]
Niech $c_n=1-2^{-n}$, $g(t)=\sum_{n\ge1}2^{-n}\ind{[c_n,1]}(t)$ na $[0,1]$ oraz $f(t)=t$. Wówczas $g$ jest niemalejąca, $V_0^1(g)=1$ i
\[
 (\KS)\int_0^1 t\,\mathrm dg(t)
 =\sum_{n=1}^\infty2^{-n}(1-2^{-n})
 =1-\frac13=\frac23.
\]
Jeśli liczymy tylko pierwszych $N$ impulsów, błąd nie przekracza $\|f\|_\infty\sum_{n>N}2^{-n}=2^{-N}$. Jest to konkretny, certyfikowany algorytm przybliżania całki.
\end{przyklad}

\begin{przyklad}[Miara Cantora: brak wzoru $g'\,\mathrm dx$]
Niech $C:[0,1]\to[0,1]$ będzie funkcją Cantora. Funkcja ta jest ciągła, niemalejąca, $C'=0$ prawie wszędzie i $C(1)-C(0)=1$, więc jej miara Stieltjesa jest osobliwa ciągła. Dla $f(t)=t$ dostajemy $(\KS)\int_0^1 t\,\mathrm dC(t)=1/2$ (z symetrii miary Cantora), a dla $f(t)=t^2$ wartość $3/8$. Istotnie, niech $\mu=\mu_C$, a $S_0(t)=t/3$, $S_1(t)=(t+2)/3$. Miara $\mu$ jest skupiona na zbiorze Cantora $K$, ponieważ $C$ jest stała w każdej luce, a $\mu$ jest bezatomowa. Każdy przedział podstawowy poziomu $n$ ma masę $2^{-n}$ z konstrukcji funkcji Cantora. Miara $\nu=\tfrac12(S_0)_*\mu+\tfrac12(S_1)_*\mu$, gdzie $(S_j)_*\mu(A)=\mu(S_j^{-1}(A))$, jest także skupiona na $K$. Każdy przedział podstawowy poziomu $n\ge1$ jest obrazem przedziału podstawowego poziomu $n-1$ przez dokładnie jedną z map $S_j$ (na poziomie zero przyjmujemy $[0,1]$); druga mapa ma obraz rozłączny z tym przedziałem. Stąd $\nu$ przypisuje mu masę $\tfrac12\,2^{-(n-1)}=2^{-n}$. Przecięcia przedziałów podstawowych z $K$, wraz z pustym zbiorem i $K$, tworzą rodzinę zamkniętą na przecięcia i generują zbiory borelowskie w $K$, ponieważ ich średnice dążą do zera i tworzą bazę topologii względnej. Z jednoznaczności skończonych miar na takiej rodzinie wynika $\nu=\mu$. Zatem dla każdej ciągłej $\psi$ zachodzi
\[
 \int\psi\,\mathrm d\mu=\frac12\int\psi(t/3)\,\mathrm d\mu(t)+
 \frac12\int\psi((t+2)/3)\,\mathrm d\mu(t),
\]
na mocy wzoru na całkowanie względem miary obrazowej. Dla $m_j=\int t^j\,\mathrm d\mu(t)$ mamy $m_0=1$. Podstawiając kolejno $\psi(t)=t$ i $\psi(t)=t^2$, otrzymujemy $3m_1=m_1+1$ oraz $9m_2=m_2+2m_1+2$. Stąd $m_1=1/2$ i $m_2=3/8$. Zastąpienie $\mathrm dC$ przez $C'(t)\,\mathrm d\lambda_1(t)$ dałoby błędnie zero. Dlatego w twierdzeniu~\ref{tw:ks-decomp} osobliwa część ciągła jest niezbędna.
\end{przyklad}

\subsection*{Całkowanie przez części ze wszystkimi poprawkami skokowymi}
\begin{twierdzenie}[Wzór całkowania przez części]
\label{tw:ks-ibp}
Niech $f,g:[a,b]\to\R$ będą regulowane i niech przynajmniej jedna z funkcji $f,g$ ma wahanie ograniczone. Wtedy obie całki KS istnieją i
\[
 \begin{aligned}
 (\KS)\int_a^b f\,\mathrm dg+(\KS)\int_a^b g\,\mathrm df
 ={}&f(b)g(b)-f(a)g(a)\\
 &+\sum_{c\in[a,b]}\bigl(\Delta^-f(c)\Delta^-g(c)
 -\Delta^+f(c)\Delta^+g(c)\bigr).
 \end{aligned}
\]
Przyjmujemy $\Delta^-f(a)=\Delta^-g(a)=0$ i $\Delta^+f(b)=\Delta^+g(b)=0$. Szereg korekcyjny jest bezwzględnie zbieżny. W szczególności, jeśli jedna z funkcji jest ciągła, cały wyraz korekcyjny znika.
\end{twierdzenie}
\begin{proof}
\emph{Krok 1: obliczenie przy jednym wspólnym skoku.} Ustalmy $c\in(a,b)$ i zapiszmy $f_-=f(c-)$, $f_0=f(c)$, $f_+=f(c+)$ oraz analogicznie $g_-,g_0,g_+$. Punktowy wkład do sumy dwóch całek wynosi $f_0(g_+-g_-)+g_0(f_+-f_-)$. Tymczasem bezpośrednie wymnożenie daje
\[
 \begin{aligned}
 f_+g_+-f_-g_-={}&f_0(g_+-g_-)+g_0(f_+-f_-)\\
 &+(f_+-f_0)(g_+-g_0)-(f_0-f_-)(g_0-g_-).
 \end{aligned}
\]
Po przeniesieniu dwóch ostatnich składników na drugą stronę otrzymujemy poprawkę
\[
 \Delta^-f(c)\Delta^-g(c)-\Delta^+f(c)\Delta^+g(c).
\] Na końcu $a$ przyjmujemy w tej tożsamości $f_-=f_0=f(a)$ oraz $g_-=g_0=g(a)$; otrzymujemy wyłącznie poprawkę $-\Delta^+f(a)\Delta^+g(a)$. Na końcu $b$ przyjmujemy $f_+=f_0=f(b)$ oraz $g_+=g_0=g(b)$; pozostaje poprawka $\Delta^-f(b)\Delta^-g(b)$.

\emph{Krok 2: funkcje schodkowe.} Wybieramy wspólny skończony zbiór węzłów obu funkcji. Na otwartych przedziałach pomiędzy nimi obie funkcje są stałe. Całki liczymy z twierdzenia~\ref{tw:ks-step}, sumując wkłady w węzłach. Dla iloczynu $p=fg$, również schodkowego, suma wszystkich przyrostów jednostronnych wynosi
\[
 \Delta^+p(a)+\sum_{c\in D\cap(a,b)}
 (\Delta^-p(c)+\Delta^+p(c))+\Delta^-p(b)
 =p(b)-p(a),
\]
gdzie $D$ jest wspólnym zbiorem węzłów. Istotnie, na każdej otwartej komórce wartość $p$ jest stała, więc przyrosty od węzła do następnej komórki i od niej do kolejnego węzła teleskopują. Zastosowanie tożsamości z kroku 1 do każdego węzła daje sumę obu całek równą $p(b)-p(a)$ powiększonej o sumę poprawek.

\emph{Krok 3: przejście do funkcji regulowanych.} Możemy bez straty ogólności założyć $f\in\BV[a,b]$, ponieważ teza jest symetryczna względem $f,g$. Istnienie całek wynika wtedy z twierdzenia~\ref{tw:ks-bv-reg} dla $\int f\,\mathrm dg$ oraz z twierdzenia~\ref{tw:ks-reg-bv} dla $\int g\,\mathrm df$. Funkcję $g$ przybliżamy jednostajnie funkcjami schodkowymi $g_n$. Dla ustalonej funkcji schodkowej $g_n$ otrzymujemy tożsamość najpierw dla schodkowych przybliżeń $f_m$ funkcji $f$ i przechodzimy do granicy przy $m\to\infty$: całka $\int f_m\,\mathrm dg_n$ jest ciągła w normie jednostajnej z oszacowaniem przez $V(g_n)$, a całka $\int g_n\,\mathrm df_m$ jest ciągła względem normy jednostajnej integratora dzięki lematowi~\ref{lem:ks-summation}, zastosowanemu do $g_n$ jako funkcji podcałkowej o wahaniu ograniczonym. Sumy poprawek zawierają tylko skończenie wiele punktów skokowych $g_n$, więc również przechodzą do granicy. Istnienie potrzebnych schodkowych przybliżeń $f_m\to f$ wynika z regulowania $f$.

Przed przejściem do granicy zauważmy, że szereg poprawek dla $f,g$ jest bezwzględnie zbieżny: $\sup_c|\Delta^\pm g(c)|\le2\|g\|_\infty$, a z dowodu twierdzenia \ref{tw:ks-bv-split} suma modułów wszystkich skoków jednostronnych $f$ (wraz z odpowiednimi końcami) nie przekracza $V_a^b(f)$. Zatem suma modułów obu rodzin iloczynów jest nie większa niż $2\|g\|_\infty V_a^b(f)$. 

Przy $n\to\infty$ całki $\int f\,\mathrm dg_n$ zbiegają do $\int f\,\mathrm dg$ z twierdzenia~\ref{tw:ks-bv-reg}; całki $\int g_n\,\mathrm df$ zbiegają dzięki lematowi~\ref{lem:ks-var}. Skoki $g_n$ zbiegają jednostajnie do odpowiednich jednostronnych skoków $g$, a sumę ich iloczynów ze skokami $f$ kontrolujemy przez $4V(f)\|g_n-g\|_\infty$: przykładowo $\sum_c|\Delta^-f(c)\Delta^-(g_n-g)(c)|\le2\|g_n-g\|_\infty V(f)$. Druga suma podlega temu samemu oszacowaniu. Wyrazy korekcyjne zbiegają więc bezwzględnie. Dowód zakończony.
\end{proof}

\begin{wniosek}[Klasyczna postać i wzór dla funkcji absolutnie ciągłej]
\label{wn:ks-ibp-ac}
Jeśli $f\in\AC[a,b]$, a $g:[a,b]\to\R$ jest regulowana, to
\[
 (\KS)\int_a^b f\,\mathrm dg
 =f(b)g(b)-f(a)g(a)-\int_a^b g(t)f'(t)\,\mathrm d\lambda_1(t).
\]
Jeśli $f\in C^1[a,b]$, ostatnią całkę można rozumieć jako zwykłą całkę Riemanna. Nie wymagamy w tym wniosku, aby $g$ miała wahanie ograniczone.
\end{wniosek}
\begin{proof}
Funkcja absolutnie ciągła ma wahanie ograniczone i jest ciągła, więc wzór z twierdzenia~\ref{tw:ks-ibp} nie ma poprawek skokowych. Ponadto $\mathrm df=f'(t)\,\mathrm d\lambda_1(t)$; z twierdzenia~\ref{tw:ks-decomp}, zastosowanego do całki $\int g\,\mathrm df$ (integrator $f\in\BV$, funkcja podcałkowa $g$ jest regulowana), otrzymujemy $\int g\,\mathrm df=\int gf'\,\mathrm d\lambda_1$. Iloczyn $gf'$ jest całkowalny w sensie Lebesgue'a, ponieważ $g$ jest ograniczona i borelowska jako jednostajna granica funkcji schodkowych, a $f'\in L^1$. Jeśli $f'\in C[a,b]$, iloczyn $gf'$ jest regulowany, ponieważ istnieją granice jednostronne obu czynników. Z przybliżenia funkcjami schodkowymi i twierdzenia \ref{ma1:thm:riemann-uniform-limit} wynika jego całkowalność w sensie Riemanna.
\end{proof}

\begin{przyklad}[Rachunek bez znajomości pochodnej funkcji całkującej]
Niech $g:[0,1]\to\R$ będzie dowolną funkcją regulowaną (nie zakładamy $g\in\BV$). Dla $f(t)=t^2$ otrzymujemy
\[
 (\KS)\int_0^1 t^2\,\mathrm dg(t)=g(1)-2\int_0^1 t\,g(t)\,\mathrm d\lambda_1(t).
\]
Ponieważ funkcja regulowana jest ograniczona i Riemanna całkowalna, ostatnia całka jest zwykłą całką Riemanna. W szczególności wzór działa dla funkcji $g$ oscylującej z nieskończonym wahaniem, o ile ma granice jednostronne.
\end{przyklad}

\begin{przyklad}[Całkująca o nieskończonym wahaniu]
Dla $t\in(0,1]$ niech $g(t)=t\sin(1/t^2)$, a $g(0)=0$. Funkcja $g$ jest ciągła, ale nie ma wahania ograniczonego: na punktach $t_n=(\pi(n+1/2))^{-1/2}$, $n\ge1$, przyjmuje wartości $(-1)^nt_n$. Dla każdego $N\ge2$ włączamy $t_N<\cdots<t_1$ do podziału przedziału $[0,1]$. Jego suma wariacyjna jest nie mniejsza niż $\sum_{n=1}^{N-1}(t_n+t_{n+1})$, co dąży do nieskończoności, ponieważ $t_n\ge (2\pi n)^{-1/2}$ dla $n\ge1$ i szereg o wykładniku $1/2$ jest rozbieżny. Pomimo tego dla $f(t)=t$ całka KS istnieje na mocy twierdzenia~\ref{tw:ks-bv-reg}, a wzór całkowania przez części daje
\[
 (\KS)\int_0^1 t\,\mathrm dg(t)
 =\sin1-\int_0^1 t\sin(1/t^2)\,\mathrm d\lambda_1(t)
 =\sin1-\frac12\int_1^\infty\frac{\sin u}{u^2}\,\mathrm du.
\]
Ostatnia całka jest bezwzględnie zbieżna i może być obliczana metodami numerycznymi. Nie potrzebujemy ani wahania ograniczonego funkcji $g$, ani klasycznej całki Riemanna--Stieltjesa.
\end{przyklad}

\begin{przyklad}[Dlaczego potrzebujemy wyrazu korekcyjnego]
Dla $H=\ind{[c,b]}$ mamy $\Delta^-H(c)=1$ i $\Delta^+H(c)=0$. Twierdzenie~\ref{tw:ks-ibp} z $f=g=H$ daje
$2(\KS)\int_a^b H\,\mathrm dH=H(b)^2-H(a)^2+1=2$. Dla $K=\ind{(c,b]}$ mamy natomiast $\Delta^-K(c)=0$, $\Delta^+K(c)=1$, zatem $2(\KS)\int_a^b K\,\mathrm dK=1-1=0$. Oba wyniki są zgodne z bezpośrednim wzorem $\int f\,\mathrm dH=f(c)$ i $\int f\,\mathrm dK=f(c)$.
\end{przyklad}

\subsection*{Całka nieoznaczona, podstawienie i błędy numeryczne}
\begin{twierdzenie}[Skoki funkcji pierwotnej KS]
\label{tw:ks-primitive}
Niech $f,g:[a,b]\to\R$ mają skończone wartości i niech
$f$ będzie KS-całkowalna względem $g$. Definiujemy
$F(x)=(\KS)\int_a^x f\,\mathrm dg$ oraz $F(a)=0$.
Jeżeli $g(c+)$ istnieje i jest skończona, to istnieje
$F(c+)$ i
\[
 F(c+)-F(c)=f(c)\Delta^+g(c).
\]
Analogicznie, jeżeli istnieje skończona granica $g(c-)$, to
\[
 F(c)-F(c-)=f(c)\Delta^-g(c).
\]
W szczególności regulowanie $g$ pociąga regulowanie $F$,
a ciągłość $g$ w $c$ pociąga ciągłość $F$ w $c$.
Jeżeli dodatkowo $f$ jest ograniczona i $g\in\BV[a,b]$, to
$F\in\BV[a,b]$ oraz $V_a^b(F)\le\|f\|_\infty V_a^b(g)$.
\end{twierdzenie}
\begin{proof}
Ustalmy $c\in[a,b)$ i $\varepsilon>0$. Wybieramy cechowanie
$\delta$ z lematu \ref{it:lem:ks-saks-henstock} dla błędu
$\varepsilon/2$. Dla $c<x\le b$ spełniającego
$x-c<\delta(c)$ pojedyncza komórka $[c,x]$ ze znacznikiem
$c$ jest drobnym podziałem częściowym. Zatem
\[
 |F(x)-F(c)-f(c)(g(x)-g(c))|\le\varepsilon.
\]
Jeżeli $g(c+)$ istnieje, to dla dostatecznie bliskiego $x>c$
mamy $|f(c)|\,|g(x)-g(c+)|<\varepsilon$ (także gdy $f(c)=0$).
Nierówność trójkąta daje
$|F(x)-F(c)-f(c)\Delta^+g(c)|<2\varepsilon$.
Dowolność $\varepsilon$ dowodzi istnienia granicy $F(c+)$
oraz wzoru na skok. Dla $c\in(a,b]$ bierzemy komórkę $[x,c]$
ze znacznikiem $c$ i otrzymujemy analogicznie
$|F(c)-F(x)-f(c)(g(c)-g(x))|\le\varepsilon$.
Przejście $x\uparrow c$ daje wzór lewostronny. Dowody
obejmują odpowiednie granice jednostronne w $a,b$.
Jeżeli $g$ jest ciągła w $c$, odpowiednie skoki są zerowe,
co oznacza ciągłość $F$ w tym punkcie.

Przy dodatkowych założeniach o $f,g$ z lokalności
(twierdzenie \ref{it:thm:ks-locality}) oraz lematu
\ref{lem:ks-var} dostajemy dla każdego podziału $P$
\[
 \sum_i|F(x_i)-F(x_{i-1})|
 \le\|f\|_\infty\sum_i V_{x_{i-1}}^{x_i}(g)
 =\|f\|_\infty V_a^b(g).
\]
Supremum po podziałach daje ostatnią tezę.
\end{proof}

\begin{twierdzenie}[Lokalna reprezentacja miarowa i dokładne wahanie funkcji pierwotnej KS]
\label{it:thm:ks-primitive-variation-exact}
Niech $g\in\BV[a,b]$, $V(t)=V_a^t(g)$, a $f:[a,b]\to\R$ będzie funkcją o skończonych wartościach, całkowalną względem uzupełnionej miary $\mu_V$. Niech
$F(t)=(\KS)\int_a^t f\,\mathrm dg$. Dla $a\le u<v\le b$ zachodzi
\[
 F(v)-F(u)=\int_{(u,v)} f\,\mathrm d\mu_g
 +f(u)\Delta^+g(u)+f(v)\Delta^-g(v).
\]
Ponadto $F\in\BV[a,b]$ i, dla $V_F(t)=V_a^t(F)$, jako równości miar borelowskich zachodzą
\[
 \mu_F=f\mu_g,\qquad \mu_{V_F}=|f|\mu_V.
\]
W szczególności $V_a^b(F)=\int_{[a,b]}|f|\,\mathrm d\mu_V$.
\end{twierdzenie}
\begin{proof}
Na mocy wniosku \ref{ma1:kh:cor:ks-variation-measure-l1}
istnieje całka KS na $[a,b]$, a z twierdzenia
\ref{it:thm:ks-locality} istnieją całki lokalne i są równe
odpowiednim przyrostom $F$.

Oznaczmy przez $\mu_g^{u,v}$ miarę przyporządkowaną
ograniczeniu $g$ do $[u,v]$; analogicznie oznaczmy
$\mu_V^{u,v}$. Z twierdzenia \ref{tw:ks-measure} wynika
\[
 \mu_g^{u,v}=\mu_g|_{(u,v)}
 +\Delta^+g(u)\delta_u+\Delta^-g(v)\delta_v .
\]
Istotnie, dla $u\le t<v$ prawa strona przypisuje
$[u,t]$ masę $g(t+)-g(u+)+g(u+)-g(u)=g(t+)-g(u)$,
gdzie dla $t=u$ pierwszy składnik jest zerowy. Dla $t=v$
masa wynosi $g(v-)-g(u+)+g(u+)-g(u)+g(v)-g(v-)$,
czyli $g(v)-g(u)$. Jednoznaczność miary daje wzór.
Identyczny argument dla $V$ oraz addytywność wahania
pokazują, że $\mu_V^{u,v}$ jest miarą wahania integratora
na $[u,v]$. Ponieważ $V$ jest niemalejąca, mamy
$\mu_V^{u,v}\le \mu_V$ (po przedłużeniu przez zero poza
$[u,v]$). Dominacja przenosi się na uzupełnienia jak w
dowodzie wniosku \ref{ma1:kh:cor:ks-variation-measure-l1}.
Możemy więc zastosować ten wniosek na $[u,v]$, co
dowodzi wzoru lokalnego.

Dla dowolnego podziału $P$ z węzłami $x_i$ mamy
\[
 \sum_i|F(x_i)-F(x_{i-1})|
 \le\sum_i\int_{[x_{i-1},x_i]}|f|\,\mathrm d\mu_V^{x_{i-1},x_i}
 =\int_{[a,b]}|f|\,\mathrm d\mu_V.
\]
Ostatnia równość wynika stąd, że wnętrza komórek są
rozłączne, a w każdym wewnętrznym węźle oba skoki
jednostronne $V$ sumują się do pełnej masy atomu.
Zatem $F\in\BV[a,b]$.

Dla $a<t<b$ ze wzoru lokalnego i twierdzenia
\ref{tw:ks-primitive} otrzymujemy
\[
 F(t+)=F(t)+f(t)\Delta^+g(t)
 =\int_{[a,t]} f\,\mathrm d\mu_g .
\]
Ta sama równość zachodzi dla $t=a$, ponieważ $F(a)=0$
i $\mu_g(\{a\})=\Delta^+g(a)$; dla $t=b$ zamiast
$F(b+)$ bierzemy $F(b)$ i stosujemy wniosek
\ref{ma1:kh:cor:ks-variation-measure-l1}.
Definicja miary $\mu_F$ i jej jednoznaczność dają
$\mu_F=f\mu_g$. Ponieważ $f$ jest całkowalna względem
$|\mu_g|$, mamy $|\mu_F|=|f||\mu_g|$: istotnie, jeśli
$s=\mathrm d\mu_g/\mathrm d|\mu_g|$, to $|s|=1$
$|\mu_g|$-prawie wszędzie, a miara $f\mu_g$ ma gęstość
$fs$. Wybieramy borelowskiego reprezentanta tej gęstości.
Miary o nieujemnych gęstościach $(fs)^+$ i $(fs)^-$
względem $|\mu_g|$ są wzajemnie osobliwe, gdyż są
skupione odpowiednio na rozłącznych zbiorach, na których
$fs>0$ i $fs<0$. Ich różnica jest równa $f\mu_g$,
zatem z jednoznaczności rozkładu Jordana ich suma jest
wahaniem $|f\mu_g|$. Ma ono gęstość $|fs|=|f|$.

Z twierdzenia \ref{tw:ks-primitive} oba skoki jednostronne
$F$ w punkcie $c$ są odpowiednimi skokami $g$
pomnożonymi przez $f(c)$. Dlatego współczynnik nadmiaru
wahania $d_c(F)$ z twierdzenia
\ref{ma1:kh:thm:ks-variation-measure-defect} jest równy
$|f(c)|d_c(g)$. Stosując to twierdzenie do $F$ i $g$,
otrzymujemy
\[
 \mu_{V_F}=|\mu_F|+\sum_c |f(c)|d_c(g)\delta_c
 =|f|\mu_V .
\]
Wzięcie masy całego przedziału daje ostatnią tezę.
\end{proof}

\begin{przyklad}[Podprzedziały rozdzielają znoszące się skoki]
\label{it:ex:ks-local-cancelled-atom}
Niech $a<c<b$, $g=\chi_{\{c\}}$ i $f=1$.
Wtedy $\Delta^-g(c)=1$, $\Delta^+g(c)=-1$,
a funkcja $g$ poza $c$ jest zerowa. Zatem $\mu_g=0$,
lecz $V$ przyjmuje wartości $0,1,2$ odpowiednio przed
$c$, w $c$ i po $c$, więc $\mu_V=2\delta_c$.
Każda suma KS z $f=1$ na $[a,t]$ teleskopuje do
$g(t)-g(a)=g(t)$; stąd $F=g$.
Wzór lokalny z twierdzenia \ref{it:thm:ks-primitive-variation-exact}
daje $(\KS)\int_a^c f\,\mathrm dg=1$ oraz
$(\KS)\int_c^b f\,\mathrm dg=-1$, mimo że całka na
$[a,b]$ jest zerowa. Ponadto $V_a^b(F)=2$.
Nie można więc zastąpić lokalnej miary integratora
prostym ograniczeniem globalnej miary $\mu_g$
ani zastąpić $\mu_V$ przez $|\mu_g|$ we wzorze na wahanie.
\end{przyklad}

\begin{twierdzenie}[Podstawienie w przedziale]
\label{tw:ks-substitution}
Niech $\varphi:[A,B]\to[a,b]$ będzie ściśle rosnącym homeomorfizmem, $\varphi(A)=a$, $\varphi(B)=b$. Niech $f,g:[a,b]\to\R$ będą funkcjami o skończonych wartościach. Wtedy całka KS po jednej stronie istnieje wtedy i tylko wtedy, gdy istnieje po drugiej, oraz
\[
 (\KS)\int_a^b f(t)\,\mathrm dg(t)
 =(\KS)\int_A^B f(\varphi(s))\,\mathrm d(g\circ\varphi)(s).
\]
Dla ściśle malejącego homeomorfizmu $\varphi:[A,B]\to[a,b]$ również zachodzi równoważność istnienia całek, przy czym
\[
 (\KS)\int_A^B f(\varphi(s))\,\mathrm d(g\circ\varphi)(s)
 =-(\KS)\int_a^b f(t)\,\mathrm dg(t).
\]
\end{twierdzenie}
\begin{proof}
Każdy podział znaczony $\mathcal Q=\{(\eta_i,[s_{i-1},s_i])\}$ przedziału $[A,B]$ przenosi się na podział $\mathcal P=\{(\varphi(\eta_i),[\varphi(s_{i-1}),\varphi(s_i)])\}$ przedziału $[a,b]$. Obie sumy KS są równe składnik po składniku. Niech $\delta$ będzie cechowaniem na $[a,b]$. Dzięki ciągłości $\varphi$ w każdym punkcie $s$ można wybrać $\widetilde\delta(s)>0$ tak, aby $|u-s|<\widetilde\delta(s)$ implikowało $|\varphi(u)-\varphi(s)|<\delta(\varphi(s))$. Wtedy każdy podział $\widetilde\delta$-drobny przenosi się na podział $\delta$-drobny, zatem istnienie i wartość całki po lewej implikują istnienie i tę samą wartość po prawej. Dla implikacji odwrotnej powtarzamy argument z $\varphi^{-1}$.
Dla malejącego homeomorfizmu odwracamy kolejność obrazów komórek, zachowując ich znaczniki. Każda komórka obrazowa ma końce w kolejności rosnącej, więc jej przyrost $g$ jest przeciwny do przyrostu $g\circ\varphi$ na komórce wyjściowej. Sumy są przeciwne; przeniesienie cechowań przez $\varphi$ i $\varphi^{-1}$ działa jak poprzednio. Stąd obie implikacje i znak minus.
\end{proof}

\begin{twierdzenie}[Podstawienie całki nieoznaczonej]
\label{tw:ks-chain}
Niech $f,g:[a,b]\to\R$ mają skończone wartości i niech
$f$ będzie KS-całkowalna względem $g$. Definiujemy
$F(x)=(\KS)\int_a^x f\,\mathrm dg$. Dla dowolnej funkcji
$h:[a,b]\to\R$ o skończonych wartościach całka KS z $h$
względem $F$ istnieje wtedy i tylko wtedy, gdy istnieje
całka KS z $hf$ względem $g$. Wówczas
\[
 (\KS)\int_a^b h(t)\,\mathrm dF(t)
 =(\KS)\int_a^b h(t)f(t)\,\mathrm dg(t).
\]
Nie wymagamy ograniczoności $h$ ani ograniczonego wahania $g$.
\end{twierdzenie}
\begin{proof}
Lokalność zapewnia istnienie $F$ na całym przedziale.
Niech $\varepsilon>0$. Rozbijamy $[a,b]$ na zbiory
$A_k=\{t:k-1\le|h(t)|<k\}$ dla $k\ge1$.
Dla każdego $k$ wybieramy cechowanie $\delta_k$ dla
całki z $f$ względem $g$ z błędem
$\varepsilon2^{-k-2}/k$.
Definiujemy $\delta(t)=\delta_k(t)$ dla $t\in A_k$.
Jeśli $\mathcal P$ jest $\delta$-drobny, to komórki
ze znacznikami w $A_k$ tworzą $\delta_k$-drobny
podział częściowy. Lemat \ref{it:lem:ks-saks-henstock}
daje dla tej grupy oszacowanie
\[
 \sum_{t_i\in A_k}|h(t_i)|
 |f(t_i)(g(x_i)-g(x_{i-1}))-F(x_i)+F(x_{i-1})|
 \le\varepsilon2^{-k-1}.
\]
Sumując po wszystkich grupach, otrzymujemy
$|S_F(h,\mathcal P)-S_g(hf,\mathcal P)|\le\varepsilon/2$.
Jeżeli jedna całka istnieje z wartością $I$, bierzemy
minimum cechowania $\delta$ i cechowania dla tej całki
z błędem $\varepsilon/2$. Nierówność trójkąta daje
błąd mniejszy od $\varepsilon$ dla drugiego rodzaju sum
względem tego samego $I$. Dowodzi to obu implikacji.
\end{proof}

\begin{twierdzenie}[Oszacowanie błędu zwykłych sum]
\label{tw:ks-quadrature}
Niech $f\in C[a,b]$, $g\in\BV[a,b]$ oraz $P=\{a=x_0<\cdots<x_n=b\}$ będzie dowolnym podziałem o średnicy $h$. Dla dowolnych znaczników $\xi_i\in[x_{i-1},x_i]$ zachodzi
\[
 \left|(\KS)\int_a^b f\,\mathrm dg-
 \sum_{i=1}^n f(\xi_i)\bigl(g(x_i)-g(x_{i-1})\bigr)\right|
 \le\omega_f(h)\,V_a^b(g),
\]
gdzie $\omega_f(h)=\sup\{|f(x)-f(y)|:x,y\in[a,b],\ |x-y|\le h\}$ jest modułem ciągłości. Jeśli $f$ jest $L$-lipschitzowska, błąd jest nie większy niż $LhV_a^b(g)$; jeśli jest h\"olderowska ze stałą $K$ i wykładnikiem $\theta\in(0,1]$, błąd nie przekracza $Kh^\theta V_a^b(g)$.
\end{twierdzenie}
\begin{proof}
Z addytywności całki KS i liniowości mamy
\[
 (\KS)\int_a^b f\,\mathrm dg-S_g(f,\mathcal P)
 =\sum_{i=1}^n(\KS)\int_{x_{i-1}}^{x_i}\bigl(f(t)-f(\xi_i)\bigr)\,\mathrm dg(t).
\]
Funkcja $f-f(\xi_i)$ ma na $[x_{i-1},x_i]$ moduł nieprzekraczający $\omega_f(h)$. Lemat~\ref{lem:ks-var} na każdym przedziale daje oszacowanie przez $\omega_f(h)V_{x_{i-1}}^{x_i}(g)$. Sumując i korzystając z addytywności wahania całkowitego, otrzymujemy żądaną nierówność. Dwa ostatnie oszacowania wynikają z definicji modułu ciągłości.
\end{proof}

\begin{uwaga}[Czego nie daje sama średnica podziału]
Jeżeli $f$ ma skok w punkcie, w którym skacze $g$, mała średnica podziału nie gwarantuje przybliżenia całki przez dowolne sumy. Dla $f=g=\ind{[c,b]}$ sumy na bardzo drobnych podziałach mogą być równe $0$ albo $1$ w zależności od znaczników. Właśnie dlatego w definicji KS stosujemy cechowania zależne od punktu. Oszacowanie z twierdzenia~\ref{tw:ks-quadrature} wymaga ciągłości $f$.
\end{uwaga}

\begin{przyklad}[Przedziałowe oszacowanie dla kwadratury]
Na $[0,1]$ niech $f(t)=\sin t$ i $g$ będzie dowolną niemalejącą funkcją o $g(0)=0$, $g(1)=3$. Dla podziału o średnicy $h=10^{-3}$ dowolna suma Stieltjesa z wartościami $\sin\xi_i$ różni się od całki KS o co najwyżej $3\cdot10^{-3}$, gdyż $|f'|\le1$. Nie potrzebujemy pochodnej $g$ ani wiedzy o położeniu jej skoków.
\end{przyklad}

\subsection*{Twierdzenia o zbieżności}
\begin{twierdzenie}[Zbieżność jednostajna przy integratorze $\BV$]\label{it:thm:ks-uniform}
Niech $g\in\BV[a,b]$ i niech $f_n$ będą KS-całkowalne
względem $g$ oraz zbiegają jednostajnie do funkcji $f$ o
skończonych wartościach. Wtedy $f$ jest KS-całkowalna
względem $g$, a całki z $f_n$ zbiegają do całki z $f$.
Funkcje $f_n$ i $f$ nie muszą być ograniczone.
\end{twierdzenie}
\begin{proof}
Od pewnego indeksu różnice $f_n-f_m$ są ograniczone, a
$\|f_n-f_m\|_\infty\to0$ przy $n,m\to\infty$.
Z liniowości oraz lematu \ref{lem:ks-var} wynika
$|I_n-I_m|\le V_a^b(g)\|f_n-f_m\|_\infty$ dla całek
$I_n$ z $f_n$. Istnieje więc skończona granica $I$.
Dla ustalonego $\varepsilon>0$ wybieramy $n$ tak duże,
aby $V_a^b(g)\|f-f_n\|_\infty<\varepsilon/3$ i
$|I_n-I|<\varepsilon/3$. Dla cechowania zapewniającego
$|S_g(f_n,\mathcal P)-I_n|<\varepsilon/3$ mamy
\[
 |S_g(f,\mathcal P)-I|
 \le V_a^b(g)\|f-f_n\|_\infty+
       |S_g(f_n,\mathcal P)-I_n|+|I_n-I|<\varepsilon.
\]
Z definicji wynika teza.
\end{proof}

\begin{twierdzenie}[Zbieżność integratorów KS]\label{ma1:kh:thm:ks-integrator-limit}
Niech $f:[a,b]\to\R$ będzie ograniczona, $g_n,g\in\BV[a,b]$ oraz $V_a^b(g_n-g)\to0$. Jeśli $f$ jest KS-całkowalna względem każdego $g_n$, to jest KS-całkowalna względem $g$ i
\[
 \left|(\KS)\int_a^b f\,\mathrm dg_n-(\KS)\int_a^b f\,\mathrm dg\right|
 \le\|f\|_\infty V_a^b(g_n-g).
\]
Jeśli natomiast $f\in\BV[a,b]$ i $g_n\to g$ jednostajnie, nie zakładając regulowania ani ograniczonego wahania $g_n$, to z istnienia wszystkich całek względem $g_n$ wynika istnienie całki względem $g$ i zbieżność całek.
\end{twierdzenie}
\begin{proof}
W pierwszym przypadku różnica dowolnych sum dla $g_n,g_m$ ma moduł nie większy niż $\|f\|_\infty V_a^b(g_n-g_m)$. Z liniowości całki i lematu \ref{lem:ks-var} oraz subaddytywności wariacji dostajemy warunek Cauchy'ego dla całek $I_n$. Niech $I$ będzie ich granicą. Dla zadanej dokładności $\varepsilon>0$ wybieramy $n$ z $|I_n-I|<\varepsilon/3$ i $\|f\|_\infty V_a^b(g_n-g)<\varepsilon/3$, a następnie cechowanie dla całki względem $g_n$ z błędem $\varepsilon/3$. Porównanie sum dla obu integratorów daje $|S_g(f,\mathcal P)-I|<\varepsilon$ dla każdego drobnego podziału. Wzór na błąd wynika z lematu \ref{lem:ks-var} zastosowanego do różnicy integratorów po wykazaniu całkowalności.

W drugim przypadku korzystamy z lematu \ref{lem:ks-summation}, którego oszacowanie sum obowiązuje dla dowolnych ograniczonych różnic integratorów:
\[
 |S_{g_n}(f,\mathcal P)-S_{g_m}(f,\mathcal P)|
 \le (2\|f\|_\infty+V_a^b(f))\|g_n-g_m\|_\infty.
\]
Zbieżność jednostajna gwarantuje ograniczoność różnic dla dostatecznie dużych indeksów. Wybierając wspólne cechowania dla dwóch całek i stosując lemat Cousina, przenosimy tę nierówność na różnicę całek. Następnie powtarzamy argument Cauchy'ego oraz wybór cechowania jak w pierwszej części, zastępując normę wariacyjną różnicy podanym oszacowaniem supremum. Nie jest potrzebne istnienie granic jednostronnych integratorów.
\end{proof}
\begin{twierdzenie}[Wspólne cechowanie i zbieżność punktowa]\label{it:thm:ks-equiintegrability}
Niech $g:[a,b]\to\R$ ma skończone wartości, $f_n$ będą
KS-całkowalne względem $g$ z całkami $I_n$, a
$f_n(t)\to f(t)\in\R$ dla każdego $t$.
Jeżeli dla każdego $\varepsilon>0$ istnieje jedno
cechowanie $\delta$ takie, że
$|S_g(f_n,\mathcal P)-I_n|<\varepsilon$ dla wszystkich $n$
i wszystkich $\delta$-drobnych podziałów $\mathcal P$,
to $f$ jest KS-całkowalna względem $g$ oraz
$(\KS)\int_a^b f\,\mathrm dg=\lim_n I_n$.
\end{twierdzenie}
\begin{proof}
Dla $\varepsilon>0$ wybieramy wspólne cechowanie z błędem
$\varepsilon/3$ i jeden drobny podział $\mathcal P$
z lematu Cousina. Jego suma ma skończenie wiele składników,
zatem $S_g(f_n,\mathcal P)\to S_g(f,\mathcal P)$.
Dla dostatecznie dużych $n,m$ różnica sum jest mniejsza
od $\varepsilon/3$. Nierówność trójkąta daje
$|I_n-I_m|<\varepsilon$. Niech $I=\lim_n I_n$.
Dla dowolnego $\eta>0$ wybieramy wspólne cechowanie
z błędem $\eta/2$. Na każdym drobnym podziale $\mathcal Q$
przechodzimy do granicy w nierówności
$|S_g(f_n,\mathcal Q)-I_n|<\eta/2$, otrzymując
$|S_g(f,\mathcal Q)-I|\le\eta/2<\eta$. To daje tezę.
\end{proof}

\begin{przyklad}[Zbieżność prawie wszędzie nie kontroluje atomów]\label{it:ex:ks-atom-convergence}
Na $[0,1]$ ustalmy $c=1/2$, $g=\chi_{[c,1]}$
oraz $f_n=\chi_{\{c\}}$ dla każdego $n$. Ciąg zbiega
do funkcji zerowej prawie wszędzie względem $\lambda_1$,
lecz $(\KS)\int_0^1 f_n\,\mathrm dg=1$ na mocy
lematu \ref{lem:ks-point}, podczas gdy całka funkcji
zerowej wynosi $0$. Zatem w twierdzeniach granicznych
nie można bez dodatkowych założeń zastąpić zbieżności
punktowej zbieżnością prawie wszędzie względem miary
Lebesgue'a. Przy ustalonej reprezentacji miarowej przez $\mu_g$ właściwą miarą dla równości prawie wszędzie jest $|\mu_g|$. Aby jednak zagwarantować istnienie całki KS dla ogólnego integratora $g\in\BV$, wystarczy kontrola względem $\mu_V$ z wniosku \ref{ma1:kh:cor:ks-measure-limits}; sama kontrola względem $|\mu_g|$ nie wystarcza (przykład \ref{ma1:kh:ex:ks-cancelled-spikes}).
\end{przyklad}

\subsection*{Granice wzorów i porównanie z całką KH}
\begin{stwierdzenie}[Nie każda całka KS jest całką Lebesgue'a--Stieltjesa]
Dla $g(t)=t$ na $[0,1]$ całka KS jest całką Kurzweila--Henstocka. Istnieją funkcje $f:[0,1]\to\R$, które są KS-całkowalne względem $g$, lecz nie należą do $L^1([0,1],\lambda_1)$. W szczególności równość z twierdzenia~\ref{tw:ks-measure} nie obowiązuje bez założeń o $f$.
\end{stwierdzenie}
\begin{proof}
Niech $F(0)=0$ i $F(t)=t^2\cos(\pi/t^2)$ dla $t\in(0,1]$ oraz $f(0)=0$, $f(t)=F'(t)=2t\cos(\pi/t^2)+(2\pi/t)\sin(\pi/t^2)$ dla $t>0$. Z twierdzenia Newtona--Leibniza dla całki KH otrzymujemy $(\KS)\int_0^1 f\,\mathrm dt=F(1)-F(0)=-1$. Drugi składnik pochodnej nie jest bezwzględnie całkowalny w otoczeniu zera. Dla $0<\eta<1$ podstawienie $u=\pi/t^2$ daje
\[
 \int_\eta^1 \frac{2\pi}{t}|\sin(\pi/t^2)|\,\mathrm dt
 =\pi\int_\pi^{\pi/\eta^2}\frac{|\sin u|}{u}\,\mathrm du.
\]
Dla każdego $n\ge1$ na przedziale $[n\pi,(n+1)\pi]$ mamy $1/u\ge1/((n+1)\pi)$, a całka z $|\sin u|$ na tym przedziale wynosi $2$. Zatem całka po tym przedziale jest co najmniej $2/((n+1)\pi)$. Rozbieżność szeregu harmonicznego dowodzi rozbieżności całki po $[\pi,\infty)$. Tymczasem pierwszy składnik należy do $L^1$. Zatem $\int_0^1|f|\,\mathrm d\lambda_1=\infty$, więc całka Lebesgue'a z $f$ jako funkcji rzeczywistej nie istnieje.
\end{proof}

\begin{uwaga}[Schemat wyboru wzoru do obliczenia]
Gdy $g$ jest schodkowa, stosujemy twierdzenie~\ref{tw:ks-step}. Gdy $g$ ma wahanie ograniczone i $f$ jest regulowana, rozkładamy miarę według twierdzenia~\ref{tw:ks-decomp}; dla funkcji całkowalnych względem $\mu_V$ stosujemy także wniosek \ref{ma1:kh:cor:ks-variation-measure-l1}. Gdy $f$ jest absolutnie ciągła, a $g$ jedynie regulowana, najlepiej zastosować wniosek~\ref{wn:ks-ibp-ac}. Przy ciągłej funkcji podcałkowej i całkującej o wahaniu ograniczonym można korzystać z oszacowania kwadraturowego~\ref{tw:ks-quadrature}. Dla dowolnego integratora i funkcji bez powyższej regularności nie istnieje jeden uniwersalny wzór redukujący całkę KS do zwykłej całki Lebesgue'a.
\end{uwaga}


\subsection{Prawostronna normalizacja integratora i atomy końcowe}
\begin{stwierdzenie}[Reprezentacja całki KS przez miarę przy prawostronnej ciągłości]\label{st:ks-miara-konwencja}
Niech $a<b$, $\alpha:[a,b]\to\R$ będzie niemalejąca,
prawostronnie ciągła na $[a,b)$ i niech $\mu_\alpha$
będzie miarą na $(a,b]$ zdefiniowaną przez $\mu_\alpha((s,t])=\alpha(t)-\alpha(s)$. Jeśli $f:[a,b]\to\R$ ma skończone wartości, a jej
ograniczenie do $(a,b]$ jest całkowalne względem uzupełnionej
miary $\mu_\alpha$ (w szczególności gdy $f$ jest regulowana),
to całka Kurzweila--Stieltjesa $\int_a^b f\,\mathrm d\alpha$ istnieje i
\[
\int_a^b f\,\mathrm d\alpha=\int_{(a,b]}f\,\mathrm d\mu_\alpha.
\]
W szczególności $\mu_\alpha(\{c\})=\alpha(c)-\alpha(c-)$ dla $a<c\le b$, a punkt $a$ nie wnosi atomu w tej konwencji.
\end{stwierdzenie}
\begin{proof}
Funkcja $\alpha$ jest monotoniczna, więc ma wahanie ograniczone. Przedłużenie $\alpha$ przez stałe poza $[a,b]$ jest niemalejące i prawostronnie ciągłe, zatem twierdzenie \atlasref{ME-04::it:thm:stieltjes-construction}{ME-04: konstrukcja miary Stieltjesa} daje miarę dodatnią, której ograniczenie do $(a,b]$ oznaczamy przez $\mu_\alpha$. Z twierdzenia \ref{tw:ks-measure} otrzymujemy miarę
kanoniczną na $[a,b]$. Jej masa w $a$ wynosi $\alpha(a+)-\alpha(a)=0$, a na przedziałach $(s,t]$ jest równa $\alpha(t)-\alpha(s)$ z prawostronnej ciągłości. Jednoznaczność skończonych miar daje więc równość miary kanonicznej z $\mu_\alpha$ przedłużoną przez zero w $a$.
Po uzupełnieniu tej miary dowolna skończona wartość
$f(a)$ nie wpływa ani na mierzalność, ani na całkowalność,
gdyż singleton $\{a\}$ ma miarę zero. Z twierdzenia
\ref{ma1:kh:thm:ks-positive-measure-l1} wynika
całkowalność KS oraz żądana reprezentacja.
Funkcja regulowana jest borelowska i ograniczona
jako jednostajna granica funkcji schodkowych, więc
spełnia przyjęty warunek dzięki skończoności miary.
Przechodząc z $s\uparrow c$ w równości $\mu_\alpha((s,c])=\alpha(c)-\alpha(s)$, otrzymujemy wskazaną masę atomu.
Dla zdegenerowanego przedziału $a=b$ oba wyrażenia
całkowe są zerowe: $(a,b]=\varnothing$, a całkę KS
na przedziale jednopunktowym definiujemy jako zero.
Nie występuje wtedy żaden skok jednostronny.
\end{proof}
\begin{twierdzenie}[Normalizacja prawostronna i atom na lewym końcu]
\label{it:thm:ks-right-normalization-endpoint}
Niech $a<b$ i $g\in\BV[a,b]$. Określmy
\[
 \widehat g(t)=
 \begin{cases}g(t+),&a\le t<b,\\g(b),&t=b.\end{cases}
\]
Wówczas $\widehat g\in\BV[a,b]$ jest prawostronnie
ciągła na $[a,b)$ i $V_a^b(\widehat g)\le V_a^b(g)$.
Jeżeli $\nu_{\widehat g}$ jest jej miarą Stieltjesa
na $(a,b]$, przedłużoną przez zero w $a$, to
kanoniczna miara KS integratora $g$ spełnia
\[
 \mu_g=\nu_{\widehat g}+
           (g(a+)-g(a))\delta_a.
\]
Dla każdej rzeczywistej funkcji regulowanej $f$ zachodzi
\[
 (\KS)\int_a^b f\,\mathrm dg
 =(\KS)\int_a^b f\,\mathrm d\widehat g
    + f(a)(g(a+)-g(a)).
\]
Jeżeli $f$ ma jedynie skończone wartości, jest całkowalna
względem uzupełnionej miary $|\mu_g|$ oraz spełnia
\[
 \sum_{c\in(a,b)}|f(c)|\,|g(c+)-g(c)|<\infty,
\]
to powyższa równość również zachodzi.
\end{twierdzenie}
\begin{proof}
Rozkład Jordana funkcji $g$ na różnicę dwóch
niemalejących funkcji pokazuje, że granice prawostronne
istnieją, a funkcja $\widehat g$ jest prawostronnie
ciągła na $[a,b)$: dla funkcji niemalejącej $p$
mamy $p(t+)=\inf_{t<s\le b}p(s)$, skąd
$\lim_{u\downarrow t}p(u+)=p(t+)$.
Dla dowolnego podziału $a=t_0<\cdots<t_m=b$
wybieramy punkty $t_i^{(n)}\in(t_i,t_{i+1})$,
$i<m$, zbieżne do $t_i$, oraz $t_m^{(n)}=b$.
Suma modułów przyrostów $g$ na tej uporządkowanej
rodzinie punktów nie przekracza $V_a^b(g)$
(dołączamy $a$ do podziału). Przechodząc do granicy
w skończonej sumie, otrzymujemy
\[
 \sum_{i=1}^m|\widehat g(t_i)-\widehat g(t_{i-1})|
 \le V_a^b(g).
\]
Supremum po podziałach dowodzi nierówności wahań.

Na mocy twierdzeń \ref{tw:ks-measure} oraz
\atlasref{ME-16::it:thm:stieltjes-signed-bv}{ME-16: miara Stieltjesa ze znakiem}, dla $a<t\le b$
mamy
\[
 \mu_g([a,t])=\widehat g(t)-g(a),\qquad
 \nu_{\widehat g}((a,t])=\widehat g(t)-\widehat g(a).
\]
Dla $t=a$ pierwsza masa wynosi
$\widehat g(a)-g(a)$, a druga miara przedłużona
przez zero nie ma masy w $a$.
Zatem obie strony żądanej równości miar zgadzają
się na wszystkich $[a,t]$. Jednoznaczność
skończonych miar ze znakiem na tej rodzinie
daje równość na zbiorach borelowskich.

Dla funkcji regulowanej obie reprezentacje
całki KS przez miarę wynikają z twierdzenia
\ref{tw:ks-measure}. Całkując względem powyższej
sumy miar, otrzymujemy wzór.
W ostatniej tezie twierdzenie
\ref{ma1:kh:thm:ks-signed-measure-l1} zapewnia
całkowalność względem $g$ i reprezentację miarową.
Ponieważ $|\nu_{\widehat g}|$ jest ograniczeniem
$|\mu_g|$ do $(a,b]$, funkcja $f$ jest również
całkowalna względem uzupełnionej miary
$|\nu_{\widehat g}|$ przedłużonej przez zero w $a$.
Dla integratora $\widehat g$ dodatkowy warunek
punktowy w tym samym twierdzeniu znika.
Otrzymujemy zatem drugą reprezentację miarową
i tę samą równość całek.
\end{proof}
\begin{uwaga}
Dla ogólnego integratora $g$ sama całkowalność
względem $|\mu_g|$ nie wystarcza do istnienia
całki KS: przeciwprzykład podano w
\ref{ma1:kh:ex:ks-cancelled-spikes}. Normalizacja
prawostronna może usuwać punktowe odchylenia
integratora istotne dla istnienia całki z funkcji
nieograniczonej, choć nie zmienia jego miary
na $(a,b]$.
\end{uwaga}

\begin{przyklad}[Skok prawostronny i wartość w punkcie]
Dla $a<c<b$ zdefiniujmy $g(x)=\chi_{(c,b]}(x)$. Funkcja $g$ nie jest prawostronnie ciągła w $c$, a jej przyrost na przedziale $[c,t]$, $t>c$, wynosi $1$. Dla dowolnej regulowanej funkcji $f$ całka KS istnieje i $\int_a^b f\,\mathrm dg=f(c)$ (jest to wkład prawego skoku). Zastąpienie $g$ przez prawostronnie ciągłą funkcję $g^+(x)=\chi_{[c,b]}(x)$ nie zmienia tu wartości całki KS dla regulowanej $f$, lecz przenosi skok z prawej strony $c$ na lewą. Zwykła miara Stieltjesa jest związana z wybranym prawostronnie ciągłym reprezentantem.
\end{przyklad}
\begin{przyklad}[Dwa skoki przeciwnego znaku w jednym punkcie]
Niech $g(x)=\chi_{\{c\}}(x)$, gdzie $a<c<b$. Funkcja ma skok lewy $\Delta^-g(c)=1$ i prawy $\Delta^+g(c)=-1$. Dla regulowanej $f$ wzór skokowy z twierdzenia \ref{tw:ks-step} daje $\int_a^b f\,\mathrm dg=f(c)-f(c)=0$. Nie należy utożsamiać takiej funkcji $g$ z dystrybuantą dodatniej miary: $g$ nie jest monotoniczna ani prawostronnie ciągła.
\end{przyklad}
\begin{przyklad}[Dwa różne końce przedziału]
Na $[a,b]$ niech $g_a=\chi_{(a,b]}$ oraz $g_b=\chi_{\{b\}}$. Dla dowolnej regulowanej $f$ zachodzi
$\int_a^b f\,\mathrm dg_a=f(a)$ oraz $\int_a^b f\,\mathrm dg_b=f(b)$ w sensie KS. Są to odpowiednio skok prawostronny w $a$ i skok lewostronny w $b$. Wzory te wynikają z twierdzenia \ref{tw:ks-step} i obowiązują także dla dowolnej $f$ o skończonych wartościach. Miara kanoniczna z twierdzenia \ref{tw:ks-measure} jest odpowiednio $\delta_a$ i $\delta_b$. Jeśli dla pierwszego integratora zmienimy wartość w $a$ na granicę prawostronną, otrzymamy funkcję stałą $1$ i miarę zerową na $(a,b]$; tracimy więc atom w $a$. Dla drugiego integratora prawostronna ciągłość na $[a,b)$ już zachodzi, a konwencja na $(a,b]$ zachowuje atom w $b$.
\end{przyklad}

\begin{przyklad}[Wartość w punkcie a wahanie miary]
\label{it:ex:stieltjes-signed-point}
Niech $0<c<1$ i $g(t)=\chi_{\{c\}}(t)$ na $[0,1]$.
Funkcja $g$ ma wahanie całkowite $2$ (na podziale
zawierającym $c$ występują przyrosty $1,-1$,
a poza nimi wszystkie przyrosty są zerowe).
Nie jest jednak prawostronnie ciągła w $c$.
Jej kanoniczna miara KS jest zerowa, ponieważ
$\Delta^-g(c)=1$, $\Delta^+g(c)=-1$,
a $g(t+)=0$ dla każdego $t<1$.
Zastąpienie $g$ prawostronnie ciągłym reprezentantem
$t\mapsto g(t+)$ (z zachowaniem wartości w $1$)
daje funkcję zerową. Wahanie tej funkcji jest zerowe
i zgadza się z wahaniem miary, jak w twierdzeniu
\atlasref{ME-16::it:thm:stieltjes-signed-bv}{ME-16: miara Stieltjesa ze znakiem}.
Różnicę między wahaniem pierwotnej funkcji $g$
i wahaniem jej miary opisuje twierdzenie
\ref{ma1:kh:thm:ks-variation-measure-defect}.
\end{przyklad}

\begin{przyklad}[Momenty miary Cantora a całki Stieltjesa]
Niech $C$ będzie funkcją Cantora, a $\mu_C$ odpowiadającą jej miarą z \atlasref{ME-17::ex:cantor-measure-17}{ME-17: miara Cantora}. Momenty tej miary obliczono w \atlasref{ME-17::it:ex:cantor-moments}{ME-17: momenty miary Cantora}.
Z twierdzenia \ref{it:thm:rs-oscillation} dla funkcji
ciągłych całki Riemanna--Stieltjesa względem $C$
są równe całkom względem $\mu_C$ (brak masy w $0$
pozwala całkować po $[0,1]$ zamiast $(0,1]$).
W szczególności
\[
\int_0^1 1\,\mathrm dC=1,\qquad \int_0^1 x\,\mathrm dC=\frac12,\qquad \int_0^1 x^2\,\mathrm dC=\frac38.
\]
Ponadto z całkowania przez części (twierdzenie \ref{it:thm:rs-tools}), ponieważ $C$ jest ciągła i ma wahanie ograniczone, wynika $2\int_0^1 C\,\mathrm dC=C(1)^2-C(0)^2=1$. Choć $C'=0$ prawie wszędzie, nie wolno zastępować $\mathrm dC$ przez $C'(x)\,\mathrm dx$: zgubilibyśmy część osobliwą miary. Wszystkie funkcje podcałkowe w tym przykładzie są
ciągłe, więc podane równości obowiązują dla całek
Riemanna--Stieltjesa, Lebesgue'a--Stieltjesa i
Kurzweila--Stieltjesa (twierdzenie \ref{tw:ks-measure}).
\end{przyklad}

\begin{przyklad}[Ta sama funkcja, różne miary]
Niech $K$ będzie zbiorem Cantora, $C$ funkcją Cantora, a $\mu_C$ miarą Cantora. Funkcja $\chi_K$ jest całkowalna w sensie Riemanna z całką $0$, ponieważ zbiór jej nieciągłości jest równy $K$ (zbiór Cantora jest domknięty i ma puste wnętrze), a $\lambda_1(K)=0$. Stosujemy twierdzenie \atlasref{ME-08::thm:riemann-criterion}{ME-08: kryterium Lebesgue'a}. Jest również całkowalna Lebesgue'a--Stieltjesa względem $\mu_C$ z całką $1$. Natomiast zwykła całka Riemanna--Stieltjesa $\int_0^1\chi_K\,\mathrm dC$ nie istnieje: punkty nieciągłości funkcji $\chi_K$ tworzą $K$, a $\mu_C(K)=1$, więc nie jest spełnione kryterium oscylacyjne z twierdzenia \ref{it:thm:rs-oscillation}.
Natomiast $C$ jest niemalejąca, a jej kanoniczna miara
KS jest równa $\mu_C$: obie mają te same wartości
na przedziałach $[0,t]$ zgodnie z twierdzeniem
\ref{tw:ks-measure}. Funkcja $\chi_K$ jest borelowska
i ograniczona, stąd całkowalna względem skończonej
miary $\mu_C$. Na mocy twierdzenia
\ref{ma1:kh:thm:ks-positive-measure-l1} otrzymujemy
\[
 (\KS)\int_0^1\chi_K\,\mathrm dC
 =\int_{[0,1]}\chi_K\,\mathrm d\mu_C=1.
\]
Ten sam integrator daje więc przykład funkcji
całkowalnej w sensie KS i Lebesgue'a--Stieltjesa,
ale nie w sensie Riemanna--Stieltjesa.
\end{przyklad}

\begin{uwaga}[Zgodność trzech konwencji]
Dla pojedynczego niezerowego skoku prawostronnie ciągłego integratora w $c\in(a,b]$ zwykła całka Riemanna--Stieltjesa wymusza ciągłość $f$ w $c$ (twierdzenie \ref{tw:rs-jeden-skok}). W całce Kurzweila--Stieltjesa cechowanie może wymusić znacznik $c$, a dla takiego skoku wkład wynosi $f(c)\Delta^-\alpha(c)$, nawet jeśli $f$ nie jest ciągła. Dla ogólnej funkcji regulowanej należy rozróżnić skok lewostronny i prawostronny: odpowiednie wkłady podano w lemacie \ref{lem:ks-point} i twierdzeniu \ref{tw:ks-bv-split}. Całka Lebesgue'a--Stieltjesa wymaga mierzalności i całkowalności względem miary, lecz nie ciągłości w atomie: wkład wynosi również $f(c)\mu_\alpha(\{c\})$. Zbieżność wartości całek nie oznacza więc równości ich dziedzin.
\end{uwaga}
\endgroup
